% concatenated from main_arxiv.tex, utils.tex
\documentclass[a4paper,11pt,fullpage]{article}

\usepackage{natbib}

\usepackage[toc,page,header]{appendix}
\usepackage{minitoc}

\usepackage{enumitem}

\usepackage{format}

\usepackage[utf8]{inputenc} % allow utf-8 input
\usepackage[T1]{fontenc}    % use 8-bit T1 fonts
\usepackage{url}            % simple URL typesetting
\usepackage{booktabs}       % professional-quality tables
\usepackage{amsfonts}       % blackboard math symbols
\usepackage{nicefrac}       % compact symbols for 1/2, etc.
\usepackage{microtype}      % microtypography
\usepackage{xcolor}         % colors
\usepackage{enumitem}       % customize list environments

\usepackage{amsmath}
\usepackage{amsthm}
\usepackage{amsfonts}
\usepackage{amssymb}
\usepackage{amsbsy}
\usepackage{mathtools}
\usepackage{graphicx}
\usepackage{color}
\usepackage{mdframed}
\usepackage{url}

\usepackage{algorithm,algorithmic}
\makeatletter
\makeatother

\makeatletter
\def\blfootnote{\gdef\@thefnmark{}\@footnotetext}
\makeatother

\usepackage[toc,page,header]{appendix}

\usepackage{enumitem}

\usepackage{tikz}
\usetikzlibrary{arrows.meta,positioning,calc}

% packages
\usepackage{algorithm,algorithmic}
\usepackage{amsmath,amsfonts,amssymb}
\usepackage{mathtools}
\usepackage{amsthm} % better to load after ams math
\usepackage{latexsym}
\usepackage{relsize}
\usepackage{enumitem}
%\usepackage{thmtools,thm-restate}
%\usepackage[amsthm]{ntheorem}

\usepackage{wrapfig}
\usepackage{float}

\newcommand{\equaldef}{\stackrel{\text{\tiny def}}{=}}


\def\elipsoid{{\mathcal E}}
\newcommand{\sym}{\mathop{\mbox{\rm sym}}}
\usepackage[capitalise]{cleveref}
\usepackage{xcolor}
\usepackage{dsfont}

\usepackage{tcolorbox}

\DeclareMathOperator*{\gain}{gain}

\newcommand{\softmax}{\text{softmax}}
\newcommand{\CH}{\text{CH}}
\def\DD{\mathcal{D}} 
\newcommand{\F}{\mathbb{F}}
\DeclareMathOperator{\vol}{volume}

\newcommand{\diag}{\mbox{diag}}

\newcommand{\uv}{\ensuremath{\mathbf u}}
\newcommand{\defeq}{\stackrel{\text{def}}{=}}


\def\ltil{{{\tilde{\ell}}}}
\def\btalpha{{\bm{\tilde{\alpha}}}}
\def\balpha{{\bm{\alpha}}}
\def\talpha{{\tilde{\alpha}}}
\def\btbeta{{\bm{\tilde{\beta}}}}
\def\bbeta{{\bm{\beta}}}
\def\tbeta{{\tilde{\beta}}}
\def\btgamma{{\bm{\tilde{\gamma}}}}
\def\bgamma{{\bm{\gamma}}}
\def\tgamma{{\tilde{\gamma}}}


\def\newregreteqn{\ensuremath{\mathrm{{AdaptiveRegret}}}}


\def\X{{\mathcal X}}
\def\H{{\mathcal H}}
\def\Y{{\mathcal Y}}
\def\W{{\mathcal W}}
\def\R{{\mathcal R}}


% notation here
\newcommand{\ftil}{\tilde{f}}

\def\hestinv{\tilde{\nabla}^{-2}}
\def\hest{\tilde{\nabla}^{2}}
\def\hessinv{\nabla^{-2}}
\def\hess{\nabla^2}
\def\grad{\nabla}


% notation

\newcommand{\re}{{\text{Re}}}
\newcommand{\B}{B}
%\newcommand{\R}{\mathcal{R}}
%\newcommand{\C}{\mathcal{C}}
\newcommand{\Q}{\mathcal{Q}}
\newcommand{\cX}{\mathcal{X}}
\newcommand{\D}{\mathcal{D}}
\newcommand{\mH}{\mathcal{H}}
\newcommand{\A}{\mathcal{A}}
\newcommand{\cE}{\mathcal{E}}
\renewcommand{\sp}{\mathrm{span}}
\newcommand{\rank}{\mathrm{rank}}
\newcommand{\cor}{\mathrm{cor}}
\newcommand{\sign}{\text{ } \mathrm{sign}}

\newcommand{\ltgain}{\mathrm{\ell_2}\text{-gain}}
\newcommand{\OPT}{\mathrm{OPT}}
\newcommand{\CR}{\mathrm{CR}}


\newcommand{\bzero}{\ensuremath{\mathbf 0}}
\newcommand{\y}{\ensuremath{\mathbf y}}
\newcommand{\z}{\ensuremath{\mathbf z}}
\newcommand{\h}{\ensuremath{\mathbf h}}
\newcommand{\K}{\ensuremath{\mathcal K}}
\def\bb{\mathbf{b}}
%\def\bc{\mathbf{c}}
\def\bB{\mathbf{B}}
\def\bC{\mathbf{C}}
\def\C{{\mathcal C}}
\newcommand{\ba}{\ensuremath{\mathbf a}}

\def\ch{\mathbf{CH}}

%\def\be{\mathbf{e}}
\def\etil{\tilde{\mathbf{e}}}
%\def\bk{\mathbf{k}}
\def\bu{\mathbf{u}}
\def\x{\mathbf{x}}
\def\calB{{\cal  B}} \def\cB{{\cal  B}}
\def\bx{\mathbf{x}}
\def\y{\mathbf{y}}
\def\w{\mathbf{w}}
\def\by{\mathbf{y}}
\def\bz{\mathbf{z}}
\def\bp{\mathbf{p}}
\def\bq{\mathbf{q}}
\def\br{\mathbf{r}}
\def\bu{\mathbf{u}}
\def\bv{\mathbf{v}}

\def\bw{\mathbf{w}}
%\def\ba{\mathbf{a}}
\def\bA{\mathbf{A}}
\def\bS{\mathbf{S}}
\def\bG{\mathbf{G}}
\def\bI{\mathbf{I}}
\def\bJ{\mathbf{J}}
\def\bP{\mathbf{P}}
\def\bQ{\mathbf{Q}}
\def\bV{\mathbf{V}}
\def\bone{\mathbf{1}}

\def\regret{\mbox{{Regret}}}

\def\ytil{\tilde{{y}}}
\def\yhat{\hat{{y}}}
\def\what{\hat{{w}}}
\def\xhat{\hat{\mathbf{x}}}
\def\xbar{\bar{\mathbf{x}}}

% comments, etc
\newcommand{\ignore}[1]{}
\newcommand{\ab}[1]{\noindent{\textcolor{blue}{\{{\bf AB:} \em #1\}}}}
\newcommand{\ip}[1]{\noindent{\textcolor{red}{\{{\bf IP:} \em #1\}}}}
\newcommand{\av}[1]{\noindent{\textcolor{magenta}{\{{\bf AV:} \em #1\}}}}

% for emails
\newcommand{\email}[1]{\texttt{#1}}

% theorems
\theoremstyle{plain}
%\theoremstyle{jmlrthm}
% \newtheorem{theorem}{Theorem}[section]
\newtheorem{theorem}{Theorem}
%\newtheorem{lemma}[theorem]{Lemma}
\newtheorem{lemma}{Lemma}
\newtheorem{corollary}[theorem]{Corollary}
%\newtheorem{proposition}[theorem]{Proposition}
\newtheorem{proposition}{Proposition}
\newtheorem{claim}[theorem]{Claim}
\newtheorem{fact}[theorem]{Fact}
\newtheorem{observation}[theorem]{Observation}
\newtheorem{assumption}{Assumption}
\newtheorem{conjecture}{Conjecture}

\newtheorem*{theorem*}{Theorem}
\newtheorem*{lemma*}{Lemma}
\newtheorem*{corollary*}{Corollary}
\newtheorem*{proposition*}{Proposition}
\newtheorem*{claim*}{Claim}
\newtheorem*{fact*}{Fact}
\newtheorem*{observation*}{Observation}
\newtheorem*{assumption*}{Assumption}

\theoremstyle{definition}
%\theoremstyle{jmlrdef}
%\newtheorem{definition}[theorem]{Definition}
\newtheorem{definition}{Definition}
%\newtheorem{remark}[theorem]{Remark}
\newtheorem{remark}{Remark}
%\newtheorem{example}[theorem]{Example}
\newtheorem{example}{Example}
\newtheorem*{definition*}{Definition}
\newtheorem*{remark*}{Remark}
\newtheorem*{example*}{Example}


% repeated theorems
 \theoremstyle{plain}
%\theoremstyle{jmlrthm}
\newtheorem*{theoremaux}{\theoremauxref}
\gdef\theoremauxref{1}

\newenvironment{repthm}[2][]{%
  \def\theoremauxref{\cref{#2}}
  \begin{theoremaux}[#1]
}{%
  \end{theoremaux}
}

%\newtheoremstyle{theoremrep}
%	{\topsep}{\topsep}              %%% space between body and thm
%	{\itshape}                      %%% Thm body font
%	{}                              %%% Indent amount (empty = no indent)
%	{\bfseries}                     %%% Thm head font
%	{}                             %%% Punctuation after thm head
%	{0.33em}                             %%% Space after thm head
%	{\cref{#3}}
%
%\theoremstyle{theoremrep}
%\newtheorem{repthm}{dontcare}



% % old refs - use \cref{} instead!
% \newcommand{\secref}[1]{Section~\ref{#1}}
%\newcommand{\secref}[1]{Sec.~\ref{#1}}
%\newcommand{\figref}[1]{Figure~\ref{#1}}
%\renewcommand{\eqref}[1]{Eq.~(\ref{#1})}
%\newcommand{\lemref}[1]{Lemma~\ref{#1}}

% \newcommand{\thmref}[1]{Theorem~\ref{#1}}
% \newcommand{\appref}[1]{Appendix~\ref{#1}}
% \newcommand{\algref}[1]{Algorithm~\ref{#1}}

% style defs
%\renewcommand{\le}{\leqslant}
%\renewcommand{\ge}{\geqslant}
%\renewcommand{\leq}{\le}
%\renewcommand{\geq}{\ge}

% bold serif
\DeclareMathAlphabet{\mathbfsf}{\encodingdefault}{\sfdefault}{bx}{n}

% operators
\DeclareMathOperator*{\err}{error}
\DeclareMathOperator*{\argmin}{arg\,min}
\DeclareMathOperator*{\argmax}{arg\,max}
\DeclareMathOperator*{\conv}{conv}
\DeclareMathOperator*{\supp}{supp}
%\DeclareMathOperator*{\poly}{poly}
\let\Pr\relax
\DeclareMathOperator{\Pr}{\mathbb{P}}

\def\mL{{\mathcal L}}
\def\mA{{\mathcal A}}

% cases
\newcommand{\mycases}[4]{{
\left\{
\begin{array}{ll}
    {#1} & {\;\text{#2}} \\[1ex]
    {#3} & {\;\text{#4}}
\end{array}
\right. }}

\newcommand{\mythreecases}[6] {{
\left\{
\begin{array}{ll}
    {#1} & {\;\text{#2}} \\[1ex]
    {#3} & {\;\text{#4}} \\[1ex]
    {#5} & {\;\text{#6}}
\end{array}
\right. }}

% macros
\newcommand{\lr}[1]{\!\left(#1\right)\!}
\newcommand{\lrbig}[1]{\big(#1\big)}
\newcommand{\lrBig}[1]{\Big(#1\Big)}
\newcommand{\lrbra}[1]{\!\left[#1\right]\!}
\newcommand{\lrnorm}[1]{\left\|#1\right\|}
\newcommand{\lrset}[1]{\left\{#1\right\}}
\newcommand{\lrabs}[1]{\left|#1\right|}
%\newcommand{\set}[1]{\{#1\}}
\newcommand{\abs}[1]{|#1|}
\newcommand{\norm}[1]{\|#1\|}
\newcommand{\ceil}[1]{\lceil #1 \rceil}
\newcommand{\floor}[1]{\lfloor #1 \rfloor}

\renewcommand{\t}[1]{\smash{\tilde{#1}}}
\newcommand{\wt}[1]{\smash{\widetilde{#1}}}
\newcommand{\wh}[1]{\smash{\widehat{#1}}}
\renewcommand{\O}{O}
%\renewcommand{\O}{\mathcal{O}}
\newcommand{\OO}[1]{\O\lr{#1}}
\newcommand{\tO}{\wt{\O}}
\newcommand{\tTheta}{\wt{\Theta}}
\newcommand{\E}{\mathbb{E}}
\newcommand{\EE}[1]{\E\lrbra{#1}}
\newcommand{\var}{\mathrm{Var}}
\newcommand{\tsum}{\sum\nolimits}
\newcommand{\trace}{\mathrm{tr}}
\newcommand{\ind}[1]{\mathbb{I}\!\lrset{#1}}
\newcommand{\non}{\nonumber}
\newcommand{\poly}{\mathrm{poly}}

% \newcommand{\tr}{^{\mathsmaller\top}}
\newcommand{\tr}{^{\mkern-1.5mu\mathsf{T}}}
\newcommand{\pinv}{^{\dagger}}
\newcommand{\st}{\star}
\renewcommand{\dag}{\dagger}

%\newcommand{\kl}[2]{\mathsf{D}_\mathrm{KL}(#1,#2)}
\newcommand{\tv}[2]{\mathsf{D}_{\mathrm{TV}}(#1,#2)}
\newcommand{\ent}{\mathsf{H}}
\newcommand{\info}{\mathsf{I}}

\newcommand{\reals}{\mathbb{R}}
\newcommand{\eps}{\varepsilon}
\newcommand{\ep}{\varepsilon}
%\newcommand{\eps}{\epsilon}
\newcommand{\sig}{\sigma}
\newcommand{\del}{\delta}
\newcommand{\Del}{\Delta}
\newcommand{\lam}{\lambda}
\newcommand{\half}{\frac{1}{2}}
\newcommand{\thalf}{\tfrac{1}{2}}

\newcommand{\eqdef}{:=}
\newcommand{\eq}{~=~}
\renewcommand{\leq}{~\le~}
\renewcommand{\geq}{~\ge~}
\newcommand{\lt}{~<~}
\newcommand{\gt}{~>~}
\newcommand{\then}{\quad\Rightarrow\quad}

\let\oldtfrac\tfrac
\renewcommand{\tfrac}[2]{\smash{\oldtfrac{#1}{#2}}}


\let\nablaold\nabla
\renewcommand{\nabla}{\nablaold\mkern-2.5mu}
%\let\cdotold\cdot
%\renewcommand{\cdot}{\mkern-2.5mu\cdotold\mkern-2.5mu}
%\newcommand{\param}{\cdotold}


%\newcommand{\diag}{\mathrm{diag}}
\newcommand{\nf}[2]{\nicefrac{#1}{#2}}
\newcommand{\vct}[1]{\bm{#1}}
\newcommand{\mat}[1]{\bm{#1}}
\newcommand{\ten}[1]{\mathcal{#1}}
\newcommand{\frob}[1]{\|#1\|_2}
\newcommand{\p}[2]{\textrm{Pr}[ #1\in #2] }
\newcommand{\pe}[2]{\mathop{\textrm{Pr}}\limits_{#1}\left[#2\right] }
%\newcommand{\one}[1]{\mathds{1}[#1]}
\newcommand{\inp}[2]{\left\langle #1,#2\right\rangle}
\newcommand{\beps}{\bm{\eps}}
\newcommand{\mX}{\mathcal{X}}
\newcommand{\mY}{\mathcal{Y}}
\newcommand{\mF}{\mathcal{F}}
\newcommand{\RE}[3]{\E{#1}{\sup_{#2}#3}}
%\newcommand{\rc}[1]{R\left(#1\right)}


\newcommand{\declarecolor}[2]{\definecolor{#1}{RGB}{#2}\expandafter\newcommand\csname #1\endcsname[1]{\textcolor{#1}{##1}}}
\declarecolor{White}{255, 255, 255}
\declarecolor{Black}{0, 0, 0}
\declarecolor{Maroon}{128, 0, 0}
\declarecolor{Coral}{255, 127, 80}
\declarecolor{Red}{182, 21, 21}
\declarecolor{LimeGreen}{75, 205, 75}
\declarecolor{DarkGreen}{0, 80, 0}
\declarecolor{Purple}{146, 42, 158}
\declarecolor{Navy}{0, 0, 128}
\declarecolor{LightBlue}{84, 101, 202}
\definecolor{mydarkblue}{rgb}{0,0.08,0.45}
\usepackage{tcolorbox}

\usepackage{xcolor}	% for color

\usepackage{makecell}

\definecolor{CardinalRed}{HTML}{C41E3A}
\definecolor{Dartmouth}{HTML}{00693E}
\definecolor{SapphireBlue}{HTML}{0F52BA}

\colorlet{MyRed}{CardinalRed}
\colorlet{MyGreen}{LimeGreen}

\colorlet{MyLightRed}{MyRed!25}
\colorlet{MyLightGreen}{MyGreen!25}

\colorlet{AlertColor}{MyRed}	% for alerts
\colorlet{BadColor}{MyRed}	% for good
\colorlet{FocusColor}{MyRed}	% for alerts
\colorlet{GoodColor}{MyGreen}	% for good
\colorlet{MacroColor}{MyRed}	% 
\colorlet{MyGreen}{LimeGreen}


\title{
\Large{Online Learning on Hidden-Convex Losses via Algorithmic Equivalence: Optimal Regret, Geometric Barrier, and Bandit Feedback}
}

\author{
  \begin{tabular}{c}
    Anas Barakat$^{1}$
  \end{tabular}
  \hfil
  \hspace*{1em}
  \begin{tabular}{c}
    Andreas Kontogiannis$^{2,5}$
  \end{tabular}
  \hfil
  \hspace*{1em}
  \begin{tabular}{c}
    Vasilis Pollatos$^{3,5}$
  \end{tabular}
  \vspace{0.2em} \\ 
  \begin{tabular}{c}
    Ioannis Panageas$^{4}$
  \end{tabular}
  \hspace*{1em}
  \begin{tabular}{c}
    Antonios Varvitsiotis$^{1,5,6}$
  \end{tabular}
}

\begin{document}
\maketitle

\blfootnote{
    Affiliations:\;\;    
    $^{1}$ Singapore University of Technology and Design \;\;
    $^{2}$ National Technical University of Athens\;\;
    $^{3}$ National and Kapodistrian University of Athens \;\;
    $^{4}$ University of California, Irvine\;\;
    $^{5}$ Archimedes, Athena Research Center, Greece\;\;
    $^{6}$ National University of Singapore, Centre for Quantum Technologies\;\;
  \vspace*{0.1em}
}
\blfootnote{
Contact: \qquad
\texttt{barakat9anas@gmail.com,
    andr.kontog@gmail.com, 
    vaspoll97@gmail.com, \\
    \hspace*{8em}
    ipanagea@ics.uci.edu,antonios@sutd.edu.sg. 
}}

\begin{abstract}
We study adversarial online learning with hidden-convex losses, i.e., nonconvex losses that become convex after a nonlinear reparameterization. 
\citet{ghai-lu-hazan22} proved that, under geometric and smoothness assumptions, online gradient descent (OGD) on such nonconvex losses approximately simulates online mirror descent (OMD) on the underlying convex losses with a suitable regularizer, yielding $\mathcal{O}(T^{2/3})$ regret. 
They left open whether the optimal $\Theta(\sqrt{T})$ regret from online convex optimization can be recovered in this hidden-convex setting. We answer this question affirmatively. More specifically, via a sharper discrete-time algorithmic equivalence argument, we prove that OGD achieves $\mathcal{O}(\sqrt{T})$ regret under the same assumptions, matching the optimal worst-case rate for adversarial online convex optimization. 
We also address another open question of \citep{ghai-lu-hazan22} by clarifying the geometry required for this algorithmic equivalence. 
We replace the diagonal-Jacobian sufficient condition with a necessary-and-sufficient Hessian compatibility condition, thereby expanding the class of admissible reparameterizations. 
We complement our tight regret bound with a lower bound showing that the Hessian compatibility assumption is essential for OGD; when it fails, we construct a smooth reparameterization and an adversarial sequence of hidden-convex losses for which OGD suffers $\Omega(T)$ regret. Finally, we extend our analysis to one-point bandit feedback and prove a $\mathcal{O}(T^{3/4})$ expected regret bound for bandit OGD with spherical smoothing, matching its classical rate on convex losses.
\end{abstract}


\section{Introduction}

Online convex optimization (OCO) provides a theoretical framework for adversarial sequential decision-making: for convex Lipschitz losses, simple first-order methods such as online gradient descent (OGD) achieve the optimal $\Theta(\sqrt{T})$ regret rate \citep{zinkevich03,hazan16oco,shalev-shwartz12oco-ftml}.
Beyond convexity, the situation is substantially more delicate. 
For arbitrary nonconvex losses, sublinear regret against the best fixed decision in hindsight is generally unattainable by simple first-order methods without additional structure or oracle access \citep{krichene-et-al15hedge-continuum,agarwal-gonen-hazan19,suggala-netrapalli20,heliou-et-al20}. For instance, it is known that no deterministic algorithm can achieve sublinear regret in online nonconvex learning in general \citep[Proposition~3]{suggala-netrapalli20}. 
This motivates the search for structured classes of nonconvex online problems that retain enough geometry to permit convex-like guarantees with simple first-order algorithms.

A prominent form of benign nonconvexity is \textit{hidden convexity} \citep{bental-teboulle96hidden-cvx,li-et-al05hidden-cvx-min,fatkhullin-he-hu25stoch-hidden-cv,levin-kileel-boumal25,gorissen-et-al26hidden-cvx}. 
This structural property appears in several modern optimization problems, including neural-network training~\citep{wang-et-al22iclr,patelsolving,ergen-pilanci25,zeger-pilanci26}, reinforcement learning~\citep{hazan-et-al19,zhang-et-al20variational,zahavy-et-al21,barakat-et-al23rlgu,barakat-et-al25rlgu,dorazio2025solving}, and nonconvex games~\citep{vlatakis2019poincare,mladenovic-et-al22hidden,sakos2023exploiting,kalogiannis2024learning,gemp-et-al25cmg,kalogiannis-et-al25zs-cmg,barakat-et-al26gumg}. In these settings, the objective is nonconvex in the algorithm’s native parameters but convex after an appropriate nonlinear change of variables. 
This structure has been exploited in deterministic and stochastic optimization to obtain global convergence guarantees for nonconvex problems \citep{chen-et-al25,fatkhullin-he-hu25stoch-hidden-cv,fatkhullin-et-al25constrained-hidden-cv,bhaskara-et-al25}. 
In the online setting, however, the learner faces a distinct challenge: the loss functions which are revealed sequentially over time may even be chosen adversarially and the goal becomes to achieve low regret against the best fixed comparator in hindsight rather than convergence to the minimizer of a fixed function.\\

\citet{ghai-lu-hazan22} initiated the study of such a structured nonconvex online learning setting through algorithmic equivalence. They showed that, under suitable geometric and smoothness assumptions on the reparameterization, OGD applied to the nonconvex losses approximately simulates OMD applied to the corresponding convex losses in the reparameterized space. 
This connection led to a $\mathcal{O}(T^{2/3})$ regret bound. Their result left open whether the slower-than-convex rate is inherent to the nonconvex online learning setting. Since the identity reparameterization recovers ordinary OCO, the best possible regret one could hope for in this class with exact gradient feedback is $\Theta(\sqrt{T})$. 
This raises the central question: 
\begin{quote}
\center
\textit{
When does OGD recover the optimal regret under exact gradient feedback for nonconvex hidden-convex losses, and does the same algorithmic equivalence principle extend to the bandit feedback setting?
}
\end{quote}
\medskip 
In this paper, we answer the exact gradient feedback question affirmatively and show that the framework extends to one-point bandit feedback. 
Specifically, we prove a $\mathcal{O}(\sqrt{T})$ regret for OGD, discuss the required compatibility conditions needed for the algorithmic equivalence, show a linear-regret lower bound without it, and sublinear expected regret guarantees in the bandit setting. We elaborate on our contributions in more details in the next section. 

\subsection{Main contributions}

We study online nonconvex learning for losses that are convex under a nonlinear reparameterization which is \textit{unknown} to the learner.  
Our results address several open questions formulated by  \citet{ghai-lu-hazan22}, refining and extending their algorithmic equivalence analysis between OGD and OMD. Our main contributions are as follows:

\begin{itemize}[leftmargin=*]
    \item \textbf{Optimal regret under exact gradient feedback.}
    In the exact gradient setting, we prove that OGD achieves $\mathcal{O}(\sqrt{T})$ regret for hidden-convex losses (Theorem~\ref{thm:main}) under the geometric compatibility and smoothness assumptions considered in \citet{ghai-lu-hazan22}. 
    This improves their $\mathcal{O}(T^{2/3})$ bound and matches the optimal regret for adversarial online \textit{convex} optimization. 
    Since classical OCO is recovered by the identity reparameterization, the rate is optimal for this problem class in general. Our analysis is based on a sharper algorithmic equivalence between OGD and online mirror descent (OMD), improving the perturbation analysis underlying the previous $\mathcal{O}(T^{2/3})$ guarantee. Specifically, our proof shows that OGD on the nonconvex losses tracks OMD in the convex parameterization with sufficiently small discretization error to preserve the classical $\sqrt{T}$ regret. 

    \item \textbf{Geometry of the reparameterization.} 
    Prior work \citep{ghai-lu-hazan22} imposed a diagonal-Jacobian condition to guarantee the Hessian compatibility needed for the OGD–OMD equivalence. We relax this structural assumption and give necessary and sufficient conditions for compatibility (Proposition~\ref{prop:matrix-field-eq-Hessian}), thereby enlarging the class of reparameterizations covered by the analysis.

    \item \textbf{Geometric barrier for OGD.}
    We show that the Hessian compatibility condition is not merely a proof artifact. 
    We prove that there exists a smooth reparameterization violating compatibility and an adversarial sequence of hidden-convex losses for which OGD incurs $\Omega(T)$ regret (Theorem~\ref{thm:linear-reg-lower-bound}).

    \item \textbf{Bandit feedback.}
    We extend our analysis to the one-point bandit setting, where the learner only observes a single loss value per round. 
    Against oblivious adversaries, we prove an expected regret bound of 
    $\mathcal{O}(T^{3/4})$ under standard additional Lipschitzness and boundedness assumptions (Theorem~\ref{thm:bandit}). 
    Notably, our result matches the best known regret rate of OGD for convex losses in the bandit setting \citep{flaxman-kalai-mcmahan05}.
\end{itemize}
Together, our results establish new regret guarantees for OGD for a structured class of \textit{online nonconvex} optimization problems, extending classical results from online \textit{convex} optimization. 

\subsection{Related Work}

\noindent\textbf{Online nonconvex optimization.}
Online convex optimization is a standard framework for adversarial sequential decision-making; see, e.g., \citet{hazan16oco,shalev-shwartz12oco-ftml,orabona19book-online-learning}. 
A separate line of work studies online learning with nonconvex losses \citep{krichene-et-al15hedge-continuum,agarwal-gonen-hazan19,suggala-netrapalli20,heliou-et-al20}. 
Sublinear global regret in such settings typically requires access to strong oracles, such as sampling oracles that allow continuous-domain exponential weights \citep{maillard-munos10,krichene-et-al15hedge-continuum}, or offline optimization oracles for the nonconvex losses \citep{agarwal-gonen-hazan19,suggala-netrapalli20}. 
Our work takes a different route: rather than assuming oracle access for general nonconvex losses, we exploit hidden convexity of the loss sequence and analyze first-order algorithms through the algorithmic equivalence between OGD in the nonconvex parameterization and OMD in the convex parameterization \citep{amid2020reparameterizing,ghai-lu-hazan22,li-et-al22implicit-bias}. A few works in the literature also consider structured nonconvex losses with special structure such as the composition of a non-increasing function with a linear function \citep{zhang-et-al15}, or weaker notions of convexity such as strictly locally quasi-convexity \citep{hazan-et-al15beyond-cvx} and weak pseudo-convexity \citep{gao-li-zhang18}.\\

\noindent\textbf{Hidden-convex optimization.}
Hidden convexity has recently been studied in several offline deterministic and stochastic optimization settings \citep{chen-et-al25,fatkhullin-he-hu25stoch-hidden-cv,fatkhullin-et-al25constrained-hidden-cv,levin-kileel-boumal25}. 
These works consider nonconvex problems that admit a convex reformulation under a nonlinear transformation, often under assumptions weaker than those used in the online algorithmic-equivalence framework. 
Their focus, however, is offline optimization: the objective is fixed, or arises as an expectation over a fixed data-generating distribution. 
By contrast, we study an adversarial online setting in which the loss functions may change arbitrarily over time, and the performance criterion is regret rather than convergence to a global optimizer. 
This distinction requires a different analysis, since the learner must compete with the best fixed decision in hindsight while only observing feedback sequentially.\\

\noindent\textbf{Online learning under bandit feedback.}
The literature on online bandit learning is extensive; we refer to \citet{lattimore-szepesvari20,lattimore24banditCO} for modern treatments. 
For bandit online convex optimization, \citet{flaxman-kalai-mcmahan05} introduced a one-point gradient estimator based on spherical smoothing, showing that gradient-based online methods can be implemented when only function values are observed. 
The bandit setting has also been studied for general Lipschitz nonconvex losses and for special nonconvex classes \citep{zhang-et-al15,gao-li-zhang18}. 
Our bandit results are different in focus: we consider hidden-convex losses and use smoothing together with the OGD--OMD equivalence to transfer regret guarantees from the convex parameterization to the original nonconvex one. 
Thus, while our algorithmic template is inspired by bandit OCO, the main challenge is controlling the interaction between bandit gradient estimation, nonlinear reparameterization, and the approximation error in algorithmic equivalence.

\section{Online Learning on Hidden-Convex Losses} 
\label{sec:OHCO}

In this section, we introduce the online \textit{hidden}-convex optimization (OHCO) problem which extends the
celebrated online convex optimization (OCO) setting \citep{shalev-shwartz12oco-ftml,hazan-et-al15beyond-cvx,orabona19book-online-learning}.\\ 

\noindent\textbf{Hidden-convex losses.} Let $\X, \Y \subset \mathbb{R}^d$ be convex compact sets. Consider a sequence of loss functions $(\ell_t)$ where for any $t \geq 1,$ $\ell_t: \X \to \mathbb{R}$ 
has the following hidden-convexity structure: 
\begin{equation}
    \label{eq:hidden-convex}
\ell_t(x) = h_t(q(x))\,, \quad \forall x \in \X\,, 
\end{equation}
where $q: \X \to \Y$ is a smooth bijective reparameterization function and $h_t: \Y \to \mathbb{R}$ is a convex function. 
In this work, we will focus on a subclass of hidden-convex functions satisfying some geometric and smoothness assumptions made in \citet{ghai-lu-hazan22}. 
We defer the statement of these assumptions on~$q$ and~$(\ell_t)$ to Sections~\ref{sec:full-info}-\ref{sec:bandit-feedback}. 
While $\ell_t$ is nonconvex in general, it is convex under the  reparameterization~$q$ and we say then that the function $\ell_t$ is \textit{hidden-convex}.  
When the reparameterization $q$ is simply the identity, the function $\ell_t$ is a (truly) convex function.\\ 

\noindent\textbf{Online regret minimization.} We are interested in the online optimization setting where a learner seeks to minimize an unknown sequence of hidden-convex functions $(\ell_t)$ by choosing a sequence
$(x_t)$ where each $x_t$ is based only on the information observed before round $t$. 
At each round $t$, a learner chooses $x_t\in \X$ using an online algorithm $\A$, then an adversary reveals a loss function $\ell_t$ which is hidden-convex (see \eqref{eq:hidden-convex}) 
and the player suffers the loss $\ell_t(x_t)$. 
The goal of the learner is to achieve a good performance with respect to the best possible fixed decision in hindsight,  minimizing the (external) regret of the algorithm $\A$ defined for any time horizon $T \geq 1$ as follows:
$$
R_T :=\sum_{t=1}^T \ell_t(x_t)-\min_{x\in \X} \sum_{t=1}^T \ell_t(x)\,.
$$
Since $q:\X\to\Y$ is bijective, the regret can equivalently be written in hidden coordinates as 
$R_T = \sum_{t=1}^T h_t(z_t) - \min_{z\in\Y}\sum_{t=1}^T h_t(z)$ where $z_t = q(x_t)\,.$ 
Unlike in the classical OCO setting, the sequence $(\ell_t)$ is not supposed to be convex. 
We recover the standard OCO setting by choosing $q$ to be the identity function in \eqref{eq:hidden-convex}.\\ 

\noindent\textbf{Information settings.} 
The reparameterization function $q$ is unknown to the learner who cannot run an online algorithm on the sequence of losses~$(h_t)$ and recover a decision in the original space~$\X$ by inverting~$q$. Preconditioning by the Jacobian of $q$ (which is unknown) is not possible either. 
We consider two different settings depending on the feedback available to the learner: 
\begin{enumerate}[label=(\roman*),leftmargin=*]
\item \textit{Exact Gradient Feedback:} In this setting, we suppose that the loss function $\ell_t$ is differentiable and the learner has access to the gradient $\nabla \ell_t(x_t)$ of the loss function $\ell_t$ at each time step of the online interaction. 
\item \textit{Bandit Feedback:} In this setting, the learner has only access to a single function value call to the loss function $\ell_t$ at each round.  
\end{enumerate}

\noindent\textbf{Online Gradient Descent.} In this work, we will focus on the most famous and classical algorithm to address this problem: Online Gradient Descent (OGD) \citep{zinkevich03}. 
In the exact gradient feedback setting, the learner performs the following update for all time steps $t$:
\begin{equation}
x_{t+1} = \Pi_{\X}(x_t - \eta \nabla \ell_t(x_t))\,,
\tag{OGD}
\label{eq:ogd-main}
\end{equation}
where $\eta$ is a positive step size and $\Pi_{\X}$ is the Euclidean projection over the convex compact set~$\X$.  
In the bandit setting, the learner will use the one-point estimate of the gradient resulting in the bandit gradient descent algorithm 
proposed by \citet{flaxman-kalai-mcmahan05}. We will discuss this case in more details in Section~\ref{sec:bandit-feedback}. 

We conclude this section by introducing some general notation that will be used throughout the paper.\\ 

\noindent\textbf{Additional notation.} 
For a nonzero integer $d$, we denote by $[d] :=\{1, \dots, d\}\,.$
For any function $f: \reals^d \rightarrow \reals^d$, we denote by $J_f(x) \in \reals^{d \times d}$ the Jacobian of~$f$ at $x$. 
Given a strictly convex and differentiable function $R:\reals^d \rightarrow \reals$, the Bregman divergence induced by $R$ is defined as
$D_R(x\| y) \eqdef R(x) - R(y) - \nabla R(y)^{\top} (x-y)~.$  
For a closed convex set $\K$ and a strictly convex regularizer $R$, we use $\Pi^R_{\K}(x) \eqdef \argmin_{y \in \K} D_R(y\| x)$ to denote the Bregman projection, 
and we use the shorthand notation $\Pi_{\K} \eqdef \Pi^{\|\cdot \|^2}_{\K}$ for the Euclidean projection.  
Given a positive-definite matrix $M \in \reals^{d \times d}$, we define the norm $\|x\|_M \eqdef \sqrt{x^{\top}Mx}$ and we denote by $M^{-\top}$ the inverse transpose matrix.   
We use the notation $B_p \eqdef \{x \in \reals^d: \|x\|_p \le 1\}$ for an $\ell_p$ ball and 
$B^{+}_p$ to denote the intersection of the $\ell_p$ ball and the positive orthant. 


\section{Regret Analysis of OGD for OHCO under Exact Gradient Feedback} 
\label{sec:full-info}

In this section, we establish a $\mathcal{O}(\sqrt{T})$ regret bound for OGD in the OHCO setting described in Section~\ref{sec:OHCO} 
under exact gradient feedback, i.e. when the learner has access to gradients of the loss functions $\ell_t$. 
We start by introducing our main assumptions in Section~\ref{subsec:full-info-as} before stating our first main result in Section~\ref{subsec:full-info-reg-bound} and providing a proof sketch in Section~\ref{subsec:full-info-proof-sketch}.  

\subsection{Assumptions}
\label{subsec:full-info-as}

We consider the same set of assumptions on the reparameterization and the loss functions as in \citet{ghai-lu-hazan22}.  

\begin{assumption}[Hessian-compatible reparameterization]
\label{as:reparam}
The map $q:\X\to\Y$ is a smooth bijection onto the convex set $\Y$ and there exists a twice continuously differentiable, strictly convex regularizer 
$R:\Y\to\mathbb R$ such that, for all $x\in\X$, with $z=q(x)$, 
$\big[\nabla^2 R(z)\big]^{-1} = J_q(x)J_q(x)^\top\,.$
\end{assumption}

Assumption~\ref{as:reparam} is a crucial assumption which restricts the class of parameterization functions $q$. 
Examples of such parameterizations include the quadratic reparameterization with the entropy regularizer,  
the exponential reparameterization with the log-barrier regularizer and the power reparameterization with 
tempered Bregman divergences \citep{amid-et-al19robust-bitempered}. 
We refer the reader to Section~3.1 in~\citet{ghai-lu-hazan22} for a more detailed discussion on these examples. 
We will elaborate more on Assumption~\ref{as:reparam} and its usefulness in Section~\ref{subsec:full-info-proof-sketch} when discussing the proof of our main results in the exact gradient feedback setting, and in Section~\ref{sec:discussion-hessian-comp-as} where we will discuss both verifiable necessary and sufficient conditions under which it is satisfied, and its necessity. 

\begin{assumption}[Regularizer and reparameterization smoothness] 
\label{as:reg-R-reparam-q-smoothness}
There exists \(G>1\) such that:
\begin{enumerate}[label=(\roman*),leftmargin=*]
    \item The map \(q\) is \(G\)-Lipschitz on \(\X\), and the 1st and 2nd derivatives of \(q^{-1}\) are bounded by \(G\). 
    \item The regularizer $R$ is $1$-strongly convex and smooth with 1st and 3rd derivatives bounded~by~$G$.
    \item For every \(z\in\Y\), the map \(w\mapsto D_R(z\|w)\) is \(G\)-Lipschitz. 
\end{enumerate}
\end{assumption}

Our last assumption asks that the loss functions $\ell_t$ have uniformly bounded gradients and that the diameter of $\X$ is bounded. 
These assumptions are also standard in OCO. 

\begin{assumption}[Boundedness of gradients and diameter]
\label{as:boundedness-grad-bregman}
There exists $G_F > 0$ s.t. $\|\nabla \ell_t(x)\|_2 \le G_F$ for all $x\in \X\,.$
Moreover, there exists~$D_1 > 0$ such that $\sup_{z,z' \in \Y} D_R(z\| z') \le D_1$\,.
\end{assumption}

\subsection{Regret Bound}
\label{subsec:full-info-reg-bound}

We are now ready to state the main result of this section.

\begin{tcolorbox}[colframe=white!, top=2pt,left=2pt,right=2pt,bottom=2pt] 
\begin{theorem}[OGD regret under full-information feedback]
\label{thm:main}
Let Assumptions~\ref{as:reparam}-\ref{as:boundedness-grad-bregman} hold and let $T \geq 1$. Setting $\eta = \sqrt{\frac{D_1}{7 G^6 G_F^3 T}}$, 
the regret of \eqref{eq:ogd-main} run on the sequence of hidden-convex losses $(\ell_t)$ defined in~\eqref{eq:hidden-convex} using stepsize $\eta$ is bounded as follows: 
$R_T \leq \sqrt{7 D_1 G^6 G_F^3 T}$.
\end{theorem}
\end{tcolorbox}

This regret bound improves over the $\mathcal{O}(T^{2/3})$ regret bound established in \citet{ghai-lu-hazan22}. 
Under Assumptions~\ref{as:reparam} to~\ref{as:boundedness-grad-bregman}, this result for \textit{hidden-convex} losses matches the standard $\mathcal{O}(\sqrt{T})$ regret bound of OGD in online \textit{convex} optimization. 

\subsection{Proof Sketch of Theorem~\ref{thm:main}}
\label{subsec:full-info-proof-sketch}

In this section, we provide a proof sketch of Theorem~\ref{thm:main}. We provide a complete proof in Section~\ref{sec:analysis}. 
The main strategy of the proof is inspired by \citet{ghai-lu-hazan22} which shows that OGD applied to non-convex functions is an approximation of online mirror descent
applied to convex functions under reparameterization. We briefly recall the intuition behind 
this idea for the convenience of the reader to motivate Assumption~\ref{as:reparam} before discussing our new key technical steps 
driving our tighter analysis.\\ 

\noindent\textbf{Implicit OMD via OGD reparameterization.} In continuous time, \citet{amid2020reparameterizing} showed that 
the mirror-flow trajectory in the hidden space and the gradient-flow trajectory in the original space coincide exactly under the reparameterization. 
More precisely, if $z(t)=q(x(t))$, they showed that both trajectories of the following ordinary differential equations: 
\begin{equation}
\label{eq:contin-flows}
\frac{d}{dt}\nabla R(z(t)) = -\nabla h(z(t))\,, 
\quad \frac{dx(t)}{dt} = -\nabla (h\circ q)(x(t))
\end{equation}
coincide if the mirror descent regularizer $R$ 
and the reparameterization function $q$ satisfy $[\nabla^2R(q(x))]^{-1} = J_q(x)J_q(x)^\top$. 
In discrete time, OGD iterates on the losses $(\ell_t)$ are given by~\eqref{eq:ogd-main} and Online Mirror Descent (OMD) iterates on the convex losses $(h_t)$ are defined as follows: 
\begin{equation}
y_{t+1} =
\arg\min_{y\in\Y}
    \left\{
        \langle \nabla h_t(y_t),y-y_t\rangle
        +
        \frac1\eta D_R(y\|y_t)
    \right\}\,,
    \tag{OMD}
    \label{eq:omd-main}
\end{equation}
with $y_0 = q(x_0)\,.$ 
In discrete time, the exact equivalence in the continuous dynamics in \eqref{eq:contin-flows} no longer holds because of the discretization error. 
Starting from a coupled state $y_t=q(x_t)$, one step of OGD on $(\ell_t)$ in the original space produces \(x_{t+1}\) and hence $z_{t+1}:=q(x_{t+1})$, while one step of OMD on $(h_t)$ in the hidden space produces~$y_{t+1}$. 
\citet{ghai-lu-hazan22} showed that these two hidden-space iterates are \(O(\eta^{3/2})\)-close.\\ 

\noindent\textbf{Our key improvement.} To obtain Theorem~\ref{thm:main}, our key contribution is to show a sharper $O(\eta^{2})$ estimate of the distance between two hidden-space iterates, which is small enough to preserve the classical $\sqrt T$ regret rate after applying a perturbed-OMD regret bound.

    \begin{lemma}[One-step OGD-OMD coupling]
    \label{prop:dist-minimizers-eta2-main}
    Suppose Assumptions~\ref{as:reparam}-\ref{as:boundedness-grad-bregman} hold. If the ghost OMD iterate $y_t$ is initialized at the image of the OGD iterate $z_t$, i.e., $y_t=z_t=q(x_t)$, $x_{t+1}$ is computed via one-step  \eqref{eq:ogd-main} and $y_{t+1}$ via one-step \eqref{eq:omd-main}, then we have  $\|z_{t+1} - y_{t+1}\|_2 \leq 6G^5G_F^3\eta^2\,.$
    \end{lemma}

The proof of Lemma~\ref{prop:dist-minimizers-eta2-main} is the main new step. Starting from the coupled state \(y_t=z_t=q(x_t)\), we first localize both
\(y_{t+1}\) and \(z_{t+1}\) in an \(O(\eta)\)-neighborhood \(\Y_t\) of \(y_t\).
On this local set, both updates can be expressed as minimizers of the same
strongly convex quadratic model, up to different higher-order perturbations:
\begin{align*}
y_{t+1}
&=
\arg\min_{y\in \Y_t}
    \left\{
        \Phi_t(y)+\varepsilon_t^{\rm OMD}(y)
    \right\}, \qquad
z_{t+1}
=
\arg\min_{y\in \Y_t}
    \left\{
        \Phi_t(y)+\varepsilon_t^{\rm OGD}(y)
    \right\},
\end{align*}
\begin{equation*}
\Phi_t(y) := \langle \nabla h_t(y_t),y-y_t\rangle
+ \frac{1}{2\eta}\|y-y_t\|_{\nabla^2R(y_t)}^2.
\end{equation*}
Here \(\varepsilon_t^{\rm OMD}\) is the Taylor remainder obtained by replacing
\(D_R(y\|y_t)\) with its second-order expansion at \(y_t\), while
\(\varepsilon_t^{\rm OGD}\) is the Taylor remainder obtained after changing
variables \(y=q(x)\) in the OGD update and expanding \(q^{-1}\) at \(y_t\).
The leading quadratic terms agree exactly because Assumption~\ref{as:reparam}
identifies the OGD Euclidean metric in \(x\)-space with the mirror metric
\(\nabla^2R(y_t)\) in \(y\)-space.

The crucial departure from \citet{ghai-lu-hazan22} is how we compare these two
perturbed minimization problems. Rather than controlling the sup-norm difference
between the perturbed objectives and then translating this into a distance between
minimizers, we compare the first-order optimality conditions directly. Since
\(\Phi_t\) is \(1/\eta\)-strongly convex and the normal cone of \(\Y_t\) is
monotone, Lemma~\ref{lem:dist-minimizers-eta-times-norm} in App.~\ref{appx:proof-main-lemma} gives
\[
    \|y_{t+1}-z_{t+1}\|
    \le
    \eta
    \left\|
        \nabla \varepsilon_t^{\rm OGD}(z_{t+1})
        -
        \nabla \varepsilon_t^{\rm OMD}(y_{t+1})
    \right\|.
\]
Since both iterates lie in an $\mathcal{O}(\eta)$-neighborhood of $y_t$, we prove thats $\|\nabla \varepsilon_t^{\rm OGD}(z_{t+1})\| = \mathcal{O}(\eta)$ and $\|\nabla \varepsilon_t^{\rm OMD}(y_{t+1})\|= \mathcal{O}(\eta)$ (in Lemma~\ref{lem:norm-grad-eps-OMD}-\ref{lem:norm-grad-eps-OGD} in App.~\ref{appx:proof-main-lemma}). Hence $\|y_{t+1}-z_{t+1}\| = \mathcal{O}(\eta^2)\,.$ 
This sharper one-step coupling improves the $\mathcal{O}(\eta^{3/2})$ discrepancy
used in \citet{ghai-lu-hazan22} to $\mathcal{O}(\eta^2)$.
Finally, the mapped OGD iterates $z_t=q(x_t)$ can be viewed as OMD iterates
perturbed by errors \(r_{t+1}=z_{t+1}-y_{t+1}\) s.t.
\(\|r_{t+1}\|\le C_{\eta} := 6G^5G_F^3\eta^2\). Applying the perturbed-OMD regret bound
(Lemma~\ref{lem2}) yields: 
\begin{equation*}
R_T \le \frac{C_{\eta} T G}{\eta} + \frac{D_1}{\eta} +\frac{\eta G^2 T}{2}\,,
\end{equation*}
which concludes the proof by optimizing the step size $\eta$. See App.~\ref{sec:analysis} for the full proof.  

\section{Discussion of the Hessian-Compatibility Hidden Geometry Assumption~\ref{as:reparam}}
\label{sec:discussion-hessian-comp-as}

Our analysis relies on Assumption~\ref{as:reparam} which allows 
to obtain an implicit OMD reparameterization and exploit the connection with OMD analysis over convex functions 
to obtain our tight regret analysis.  
In this section, we discuss this assumption in more depth. First, we answer an open question 
in \citet{ghai-lu-hazan22} regarding whether we can relax their sufficient diagonal Jacobian condition (Assumption 4 therein) 
to verify Assumption~\ref{as:reparam}. In section~\ref{subsec:relaxing-diag-jacobian}, we provide necessary and sufficient conditions under which Assumption~\ref{as:reparam} is satisfied. 
Second, we address the question of the necessity of this Hessian-compatibility hidden-geometry assumption, which is also mentioned as an open question in 
\citet{ghai-lu-hazan22}. In Section~\ref{subsec:lin-OGD-regret-without-Hessian-comp}, we show that Assumption~\ref{as:reparam} is necessary to obtain 
sublinear regret for OGD. 

\subsection{Relaxation of the diagonal Jacobian assumption in \citet{ghai-lu-hazan22}}
\label{subsec:relaxing-diag-jacobian}

Define the inverse metric field $M:\Y\to\mathbb R^{d\times d}$ by
\begin{equation}
\label{eq:matrix-field-M}
    M(z) := \left[ J_q(q^{-1}(z))J_q(q^{-1}(z))^\top \right]^{-1}.
\end{equation}
Recall that Assumption~\ref{as:reparam} states that there exists a scalar function $R:\Y\to\mathbb R$ s.t. $M(z)=\nabla^2 R(z),$ for all $z\in\Y.$ 
Theorem 7 in \citet{ghai-lu-hazan22} showed that if $J_q$ is diagonal then there exists such $R$. They used this diagonal separable structure  
to construct a regularizer $R$ explicitly as a solution to second-order ODEs with coefficients depending on the parameterization function $q$ and its derivatives. 
In the following result, we relax this condition and provide more general conditions under which the matrix field $M$ is a Hessian field. 

\begin{tcolorbox}[colframe=white!, top=2pt,left=2pt,right=2pt,bottom=2pt] 
\begin{proposition}
\label{prop:matrix-field-eq-Hessian}
If $q$ is a $C^3$ diffeomorphism and the matrix field $M$ defined in \eqref{eq:matrix-field-M} satisfies:  
\begin{equation*}
\partial_{z_k} M_{ij}(z) = \partial_{z_j} M_{ik}(z)\,,\quad \forall i, j, k \in [d],\quad  \forall z \in \Y\,,
\end{equation*}
where $\Y$ is a convex domain, then there exists a scalar function $R$ s.t. $M = \nabla^2 R$, which is unique up to an affine function. Moreover, the function $R$ is also strongly convex on $\Y$. 
\end{proposition}
\end{tcolorbox}

This result follows as a consequence of the fundamental result from differential analysis stating that a vector field with 
matching cross-partial coordinate function derivatives is a conservative field, i.e., a gradient field deriving from a scalar potential.  
See Proposition~\ref{prop:existence-grad-potential} for a precise statement of that result and a complete proof of Proposition~\ref{prop:matrix-field-eq-Hessian}. 

We provide two examples where the Jacobian of the parameterization $q$ is not diagonal in general, yet the matrix field $M$ defined in \eqref{eq:matrix-field-M} satisfies the cross-derivatives conditions $\partial_{z_k}M_{ij}=\partial_{z_j}M_{ik}$ of Proposition~\ref{prop:matrix-field-eq-Hessian}.  
Moreover, the compatibility Assumption~\ref{as:reparam} is satisfied with explicit regularizers $R$. 
Detailed derivations for both examples can be found in App.~\ref{appx:detailed-examples}. 

\begin{example}[Affine mixing of a separable parameterization]
Let $q(x)=As(x)+b$ and $s(x)=\bigl(s_1(x_1),\dots,s_d(x_d)\bigr)$ where \(A\in\mathbb R^{d\times d}\) is invertible, \(b\in\mathbb R^d\), and each
\(s_i:\mathbb R\to\mathbb R\) is a \(C^2\) diffeomorphism. Then $J_q(x)=A\,D(x)$ with $D(x):=\operatorname{diag}\bigl(s_1'(x_1),\dots,s_d'(x_d)\bigr)$. 
Hence $M(z) = A^{-T}D(x)^{-2}A^{-1}\,,z=q(x).$ Moreover, $M$ is explicitly the Hessian of the following regularizer: $R(z):=\sum_{i=1}^d r_i\bigl((A^{-1}(z-b))_i\bigr)$ and $r_i''(t) = \bigl(s_i'(s_i^{-1}(t))\bigr)^{-2}\,.$
For instance, if $s_i(u_i) = \exp(u_i)$ for all $i \in [d],$ then $r_i(t) = - \log t$ satisfies the conditions. 
\end{example}

\begin{example}[Rank-one nonlinear mixing]
Let \(a\in\mathbb R^d\) and let $q(x)=x+h(a^\top x)\,a$ 
where \(h:\mathbb R\to\mathbb R\) is \(C^2\). We have $J_q(x)=I+h'(a^\top x)\,aa^\top.$
This Jacobian is typically non-diagonal whenever at least two components of $a$ are nonzero.
\end{example}

\subsection{OGD can incur linear regret without Hessian-compatible hidden geometry}
\label{subsec:lin-OGD-regret-without-Hessian-comp}

 In this section, we discuss how crucial is Assumption~\ref{as:reparam} to obtain sublinear regret for OGD 
in our nonconvex online learning setting. Perhaps surprisingly, we show that OGD can incur linear regret 
if the Hessian compatibility assumption does not hold; i.e., if the metric induced by the symmetric positive definite matrix $M$ 
in \eqref{eq:matrix-field-M} does not derive from a Hessian metric such that $M(z) = \nabla^2 R(z)$ for some regularization function $R$. 

\begin{tcolorbox}[colframe=white!, top=2pt,left=2pt,right=2pt,bottom=2pt] 
\begin{theorem}[Linear regret without Hessian-compatible geometry]
\label{thm:linear-reg-lower-bound}
There exist a smooth reparameterization $q$ with uniformly positive definite induced metric $J_q(x)J_q(x)^\top \succeq \mu I$ whose 
inverse metric field is not the Hessian of any regularizer, and an adaptive sequence of time-varying convex hidden losses $h_t$ 
such that OGD run on the original losses $\ell_t = h_t \circ q$ suffers linear regret: $R_T = \Omega(T)$. 
\end{theorem}
\end{tcolorbox}

The theorem shows that smoothness and uniform positive definiteness of the induced hidden geometry are not sufficient for sublinear regret of OGD. This means that Assumption~\ref{as:reparam} is not merely a technical condition for our analysis of OGD in Theorem~\ref{thm:main}. 
The obstruction is geometric: when the metric induced by the reparameterization is not Hessian-compatible, OGD may fail to achieve sublinear regret even when the reparameterization is smooth and its induced metric is uniformly positive definite.
The hidden-space dynamics of OGD need not arise from an integrable potential. 
An adversary can then exploit this fact by forcing the iterates to cycle around a loop with nonzero circulation and force linear regret. We elaborate more on this in the next proof sketch.\\ 

\noindent\textbf{Proof sketch of Theorem~\ref{thm:linear-reg-lower-bound}.} We provide a brief proof sketch of the result to highlight the intuitions. An extended proof sketch with all the steps can be found in App.~\ref{app:detailed-proof-sketch-thm-linear-reg-lower-bound} and the full proof is provided in App.~\ref{appx:complete-proof-thm:linear-reg-lower-bound}. 
The construction uses the two-dimensional reparameterization
$q(x_1,x_2)=e^{x_1}(\cos x_2,\sin x_2)$ on a compact rectangular domain. 
It can be easily verified that its induced metric is uniformly positive definite, but the corresponding inverse metric in hidden coordinates is $M(z)= [J_q(x) J_q(x)^{\top}]^{-1} = \frac{1}{\|z\|^2}I_2$ where $z = q(x)$, which is not Hessian-compatible in the sense needed for the OGD-OMD equivalence.
We consider linear hidden losses $\ell_t(x)=h_t(q(x))$ with $h_t(z)=\langle s_t,z\rangle$. 
Writing $z_t=q(x_t)$, the OGD update in the original coordinates induces the hidden-space dynamics: 
\begin{equation}
\label{eq:ogd-dyn-hidden-space-main}
    z_{t+1} = z_t-\eta M(z_t)^{-1}s_t + O(\eta^2).
\end{equation}
The adversary chooses $s_t=-M(z_t)v_t$, where $v_t$ is a unit coordinate direction. 
Thus the OGD iterates in the hidden space satisfy:
\begin{equation*}
    z_{t+1}=z_t+\eta v_t+O(\eta^2),
\end{equation*}
so the adversary can make the trajectory approximately follow any axis-aligned polygonal path.
Fixing a comparator $u \in \X$ and writing $z_u=q(u) \in \Y$, the instantaneous regret can be expressed~as: 
\begin{equation*}
    \ell_t(x_t)-\ell_t(u)
    =
    -\langle F_u(z_t),v_t\rangle,
    \qquad
    F_u(z):=M(z)(z-z_u).
\end{equation*}
Using the hidden-space dynamics~\eqref{eq:ogd-dyn-hidden-space-main}, it follows that: 
\begin{equation*}
    \ell_t(x_t)-\ell_t(u)
    =
    -\frac{1}{\eta}
    \langle F_u(z_t),z_{t+1}-z_t\rangle
    +O(\eta).
\end{equation*}
Hence cumulative regret along a trajectory is, up to lower-order terms, a rescaled discrete line integral of the vector field $F_u$. For an appropriate comparator $u$, the vector field $F_u$ is not conservative. Equivalently, its curl is negative on a small rectangle $R$ in hidden space, so the counter-clockwise circulation satisfies $\oint_{\partial R} F_u(z)\cdot dz < 0.$ 
The adversary chooses the directions $v_t\in\{\pm e_1,\pm e_2\}$ (where $e_1 = (1,0)$ and $e_2 = (0,1)$) so that the hidden iterates move once around this rectangle. 
The Riemann-sum approximation of the line integral then gives regret of order $1/\eta$ over one cycle, while the cycle itself lasts $\Theta(1/\eta)$ rounds. 
Thus each cycle contributes a positive constant amount of regret per round. 
Repeating the same cycle for the horizon $T$ yields $R_T=\Omega(T)$. 
The full construction and error estimates are given in App.~\ref{app:proof-thm-linear-reg-lower-bound}. 

\begin{remark}
The proof of Theorem~\ref{thm:linear-reg-lower-bound} goes through for any reparameterization function $q$ such that $[J_qJ_q^{\top}]^{-1}$ is not the Hessian of a function $R$ and the vector field $F_u = (P_u, Q_u)$ defined in \eqref{eq:def-F_u-main} satisfies $\partial_{z_1} Q_u(z) - \partial_{z_2} P_u(z) < 0$ (hence is not conservative). This set of reparameterization functions is not restricted to the specific function $q$ used in the lower bound construction.   
\end{remark}

\noindent\textbf{Open question.} The lower bound of Theorem~\ref{thm:linear-reg-lower-bound} only holds for OGD. 
Whether Assumption~\ref{as:reparam} is necessary to obtain sublinear regret in general, for \textit{any} online algorithm, remains open. 


\section{Regret Analysis of OGD for OHCO under Bandit Feedback} 
\label{sec:bandit-feedback}

\subsection{Expected Regret Bound in the Bandit Setting}

We now consider the one-point bandit feedback, where the learner observes only one function value per round. 
We analyze the classical Bandit OGD (BOGD) algorithm with spherical smoothing (Algorithm~\ref{alg4}, \citet{flaxman-kalai-mcmahan05}) in the hidden-convex setting, against an oblivious adversary. 
Let $U_{sp}$ denote the uniform distribution over the Euclidean unit sphere. 
For an exploration radius $\delta>0$, the algorithm maintains its iterates in the shrunken domain $\X_\delta := (1-\delta)\X$, so that the perturbed query points remain feasible (under the standard normalization $0\in\X$ and $B_2\subseteq\X$ as in \citet{flaxman-kalai-mcmahan05}). 
We denote by $\Y_\delta := q(\X_\delta)$ the corresponding hidden-space domain, which is only used in the analysis to compare the mapped BOGD iterates with mirror descent on the convex losses.\\

\begin{wrapfigure}{r}{0.45\linewidth}
\vspace{-2\baselineskip}
\begin{minipage}{\linewidth}
\begin{algorithm}[H]
\caption{Bandit OGD (BOGD) \citep{flaxman-kalai-mcmahan05}} \label{alg4}
\begin{algorithmic}[1]
\STATE Input: $x_1 \in \X_\delta$\\
\FOR{$t = 1, \ldots, T$}
\STATE Sample a random direction $\zeta_t \sim U_{sp}$
\STATE Choose $\hat{x}_t = {\color{blue} x_t} + \delta \zeta_t$
\STATE Observe the bandit loss $\ell_t(\hat{x}_t)$
\STATE Construct ${\color{MyGreen} g_t} = \frac{d}{\delta} \ell_t(\hat{x}_t) \zeta_t$
\STATE 
Update 
$x_{t+1} = \Pi_{\X_{\delta}} ({\color{blue} x_t}-\eta {\color{MyGreen} g_t})$
\ENDFOR
\end{algorithmic}
\end{algorithm}
\end{minipage}
\vspace{-3\baselineskip} 
\end{wrapfigure}

\noindent Before stating the result, we introduce a few extra standard assumptions on the smoothness and boundedness of the losses $(\ell_t)$, and the diameter of~$\X$. 

\begin{assumption}[Loss boundedness]
There exists $M>0$ s.t. $|\ell_t(x)| \le M$, for all $x \in \X, t \geq 1$. 
\end{assumption}
\begin{assumption}[Smoothness]
\label{as:smoothness-bandit}
    For any $t$, the loss function $\ell_t$ is smooth, i.e., 
    there exists a constant $H>0$, such that for any $x,x' \in \X$ and any $t \geq 1$, 
    $\|\nabla \ell_t(x) - \nabla \ell_t(x')\| \leq H \|x - x'\|.$
\end{assumption}

\begin{assumption}[Boundedness of $\X$]
\label{as:bounded-euclidean-diam}
    There exists a constant $D>0$, such that for  any $x \in \X$, it holds that $\|x \|_{2} \le D$, and for any $x_1,x_2 \in \X$, it holds that $\|x_1-x_2 \|_{2} \le D$.
\end{assumption}

In the following theorem, we present the no-regret guarantees of Bandit OGD.

\begin{tcolorbox}[colframe=white!, top=2pt,left=2pt,right=2pt,bottom=2pt] 
\begin{theorem}[Expected regret under bandit feedback]
\label{thm:bandit}
Let Assumptions~\ref{as:reparam} to~\ref{as:bounded-euclidean-diam} hold. 
Then, the expected regret of BOGD (Algorithm \ref{alg4}) is bounded as follows:
\begin{equation*}
\mathbb{E}[\text{R}_T] \leq \sqrt{dM G^3 (G_FD + G^2 HD + G_F) (3 D_1^{1/2} + D_1^{3/2})}\,\cdot T^{3/4}\,,
\end{equation*}
when using the stepsize $\eta = \left(\frac{4 D_1^3}{K_1^2 K_2 T^3}\right)^{1/4}$, and the exploration radius $\delta = \left(\frac{4 K_2 D_1}{K_1^2 T}\right)^{1/4}$, where $K_1 := G_FD + G^2 HD + G_F$ and $K_2 := 6 d^2 G^6 M^2$. 
\end{theorem}
\end{tcolorbox}

Our result implies that, despite hidden convexity, OGD with spherical smoothing achieves the same best known regret rate as in the convex bandit case \citep{flaxman-kalai-mcmahan05}. 

\subsection{Proof Sketch of Theorem~\ref{thm:bandit}}
\label{subsec:bandit-proof-sketch}

We give a proof sketch of Theorem~\ref{thm:bandit}; the complete proof is in
App.~\ref{appx:proofs-bandit}. The proof follows the same high-level strategy as in the exact gradient setting with a few similar steps to the analysis of OBGD in OCO: after mapping the BOGD iterates to the hidden
space, $z_t=q(x_t)$, we compare them to mirror-descent steps on the convex losses
$h_t$. The bandit setting introduces three additional errors: the domain must be
shrunk to allow perturbations, the one-point estimator is biased for the true
gradient, and the algorithm suffers the loss at the perturbed point
$\hat x_t=x_t+\delta\zeta_t$, rather than at $x_t$. To define the ghost OMD iterate in the hidden space (which is only used in the analysis and never computed), we introduce the hidden-space gradient estimator $\tilde g_t := J_q(x_t)^{-\top} g_t$ which is the natural hidden-space analogue of the bandit estimator $g_t$ by the chain rule: $\nabla h_t(z_t)=J_q(x_t)^{-\top}\nabla \ell_t(x_t), z_t=q(x_t).$
Note that it is not computed by the algorithm but introduced only to analyze the mapped BOGD trajectory as a perturbed OMD trajectory in $\Y_\delta=q(\X_\delta)$. Since \(|\ell_t|\le M\) and \(\|\zeta_t\|=1\), we have $\|g_t\| \leq dM/\delta$. 
Let $z^\star\in\arg\min_{z\in\Y}\sum_{t=1}^T h_t(z)$ and $z^\star_\delta\in\arg\min_{z\in\Y_\delta}\sum_{t=1}^T h_t(z)$ be two comparators, and let \(\hat z_t=q(\hat x_t)\). 
The starting point is the following regret decomposition:
\begin{multline*}
\mathbb{E}[R_T]
\le
\underbrace{
\mathbb{E}\left[\sum_{t=1}^T
\langle \tilde g_t,z_t-z^\star_\delta\rangle\right]
}_{\text{\makecell{$R_T^{(1)}:$ perturbed OMD regret\\ for random linear losses}}}
+
\underbrace{
\sum_{t=1}^T h_t(z^\star_\delta)-h_t(z^\star)
}_{\text{\makecell{$R_T^{(2)}:$ domain\\ shrinkage error}}}
+
\underbrace{
\mathbb{E}\left[\sum_{t=1}^T
B_t\right]
}_{\text{\makecell{$R_T^{(3)}$: gradient\\ estimation bias}}}
+
\underbrace{
\mathbb{E}\left[\sum_{t=1}^T h_t(\hat z_t)-h_t(z_t)\right]
}_{R_T^{(4)}:\ \text{smoothing error}},
\end{multline*}
where $B_t := \langle b_t,z^\star_\delta-z_t\rangle$ and $b_t:=\mathbb{E}[\tilde g_t\mid\mathcal F_{t-1}]-\nabla h_t(z_t)\,.$ 
The four terms isolate the distinct sources of error in the bandit setting. The shrinkage term is controlled by comparing optimization over $\X$ and
$\X_\delta=(1-\delta)\X$, giving $R_T^{(2)}\le \delta LDT$. 
The smoothing term is controlled directly by Lipschitzness of $\ell_t$: $|R_T^{(4)}| \leq \delta LT.$ 
The bias term comes from the fact that the one-point spherical estimator is unbiased for the gradient of the ball-smoothed loss, not for $\nabla\ell_t(x_t)$. Using
smoothness of \(\ell_t\) and the chain-rule transformation above yields $\|b_t\|\le \delta GH$ and $\|z_t-z^\star_\delta\|\le GD$, and hence $|R_T^{(3)}|\le \delta G^2HDT\,.$ 
It remains to control \(R_T^{(1)}\), which is the term most closely related to the exact gradient setting proof. For each round, define the ghost OMD iterate, only used in the analysis: 
\begin{equation*}
    y_{t+1}
    =
    \arg\min_{y\in\Y_\delta}
    \left\{
        \langle \tilde g_t,y-z_t\rangle
        +
        \frac1\eta D_R(y\|z_t)
    \right\}.
\end{equation*}
The actual
mapped BOGD iterate satisfies $z_{t+1}=q(x_{t+1})=y_{t+1}+r_{t+1}$. 
As in the exact gradient feedback setting, the OGD--OMD one-step coupling applies pathwise because the proof only requires a bounded vector $g_t$, not an exact gradient.
With $\hat G_F=dM/\delta$, Corollary~\ref{cor:approx-lemma-bandit} gives $\|r_{t+1}\| \leq \hat C_\eta:= G^5\left(5\hat G_F^2+\hat G_F^3\eta\right)\eta^2 .$
Similarly to the exact setting, the perturbed OMD bound for the random linear losses
\(\langle \tilde g_t,\cdot\rangle\) then gives $R_T^{(1)}
    \le
    \frac{D_1}{\eta}
    +
    \frac{\eta T}{2} \hat{G}_F^2
    + \frac{G \hat C_\eta T}{\eta}\,,$ 
    where $\hat{G}_F^2$ and $\hat C_\eta$ replace $G_F$ and $C_\eta$. 
Combining the four bounds yields: 
$\mathbb{E}[R_T]
    \le
    K_1\delta T
    +
    \frac{D_1}{\eta}
    +
    K_2\frac{\eta T}{\delta^2}
    +
    K_3\frac{\eta^2T}{\delta^3}\,.$
Balancing the first three terms gives the values of $\eta, \delta$ reported in Theorem~\ref{thm:bandit} and the desired
$\mathcal{O}(T^{3/4})$ bound follows.

\section{Conclusion and Future Work}

In this paper, we studied online nonconvex learning for a structured class of hidden-convex losses. 
We showed that OGD achieves $\mathcal{O}(\sqrt{T})$ regret in the exact-gradient setting, sharpened the geometric conditions required for the underlying algorithmic equivalence, and proved that OGD can suffer linear regret when these conditions fail. 
We also extended the framework to one-point bandit feedback, establishing sublinear expected regret guarantees against oblivious adversaries.

Several questions remain open. 
First, our lower bound shows that OGD can suffer linear regret without Hessian compatibility, but it does not rule out sublinear regret for other algorithms. 
An important direction is therefore to determine whether online hidden-convex optimization admits sublinear, or even $\mathcal{O}(\sqrt{T})$, regret without this geometric assumption. 
Second, our bandit guarantees leave open the possibility of sharper rates, high-probability bounds, and algorithms matching the best-known guarantees from bandit online convex optimization. 
More broadly, it would be interesting to understand for which classes of hidden-convex functions online algorithms incur a regret that matches that of their convex counterparts.


\section*{Acknowledgements}

Ioannis Panageas is supported by NSF grant CCF- 2454115. This work is supported by the MOE Tier 2 Grant (MOE-T2EP20223-0018),  the CQT++ Core Research Funding Grant (SUTD) (RS-NRCQT-00002), the National Research Foundation Singapore and DSO National Laboratories under the AI Singapore Programme (Award Number: AISG2-RP-2020-016), and partially by Project MIS 5154714 of the National Recovery and Resilience Plan, Greece 2.0, funded by the European Union under the NextGenerationEU Program.

\bibliographystyle{plainnat}
\bibliography{references}

\newpage
\appendix

\tableofcontents

\section{Proofs for Section~\ref{sec:full-info}: Proof of Theorem~\ref{thm:main}}
\label{sec:analysis}


In this section, we provide a complete proof of Theorem~\ref{thm:main}. The proof strategy consist in showing that the OGD iterates mapped to $\Y$ via $q$ and the OMD iterates are close to each other in the hidden space~$\Y$ after a single update step starting from the same initial point. Then we can view the OGD update as a perturbed version of OMD, and combine it with the fact that the OMD algorithm can tolerate bounded noise per trial. 

\subsection{Proof of Lemma~\ref{prop:dist-minimizers-eta2-main}}
\label{appx:proof-main-lemma}

We begin with the following key lemma showing that the updates $y_{t+1}$ and $z_{t+1} = q(x_{t+1})$ created by OMD and OGD respectively, are close to each other when starting from the same initial point $y_t= z_t = q(x_t)$. We show that the two iterates are minimizers of approximately the same strongly-convex objective, hence the minimizers must be close. 
In preparation of the bandit setting, we prove a general result that holds pathwise when we use any bounded vector instead of the exact gradients of both the original loss function $\ell_t$ and the hidden function $h_t.$ 

\begin{tcolorbox}[colframe=white!, top=2pt,left=2pt,right=2pt,bottom=2pt] 
\begin{proposition}\label{prop:dist-minimizers-eta2}
Suppose Assumptions~\ref{as:reparam}-\ref{as:boundedness-grad-bregman} hold and assume $y_t= q(x_t)$.
Consider the following OMD and OGD-like update rules: 
\begin{align}
y_{t+1}&=\argmin_{y\in \Y} \tilde{g}_t^{\top}(y-y_t) +\frac{1}{\eta} D_R(y\|y_t)\,, \label{eq:omd}\\
x_{t+1}&=\argmin_{x\in \X} g_t^{\top}(x-x_t) +\frac{1}{2\eta} \|x-x_t\|_2^2 \label{eq:ogd}\,, 
\end{align}
where $\tilde{g}_t = J_q(x_t)^{-\top} g_t$ and $g_t \in \mathbb{R}^d$ is any vector\footnote{$\tilde{g}_t$ is used in the analysis but never actually computed, in contrast to $g_t$ which is computed and used in BOGD. Note that we suppose that we do not have access to the parameterization $q$.}. If in addition there exists $\hat{G}_F > 0$ such that $\|g_t\| \leq \hat{G}_F$ for any $t$, then,  
\begin{equation}
\label{eq:approx-ineq}
    \|y_{t+1}-q(x_{t+1})\|_2 \leq G^5 (5\hat{G}_F^2 + \hat{G}_F^3 \eta) \eta^2\,.
\end{equation} 
\end{proposition}
\end{tcolorbox}

Before providing the proof, we state an immediate corollary of this result in the exact gradient information setting of Theorem~\ref{thm:main}. 

\begin{tcolorbox}[colframe=white!, top=2pt,left=2pt,right=2pt,bottom=2pt] 
\begin{corollary}[Proposition~\ref{prop:dist-minimizers-eta2} in the exact gradient setting]
\label{cor:approx-lemma-det}
In Proposition~\ref{prop:dist-minimizers-eta2}, set $g_t = \nabla \ell_t(x_t)$ and $\tilde{g}_t = \nabla h_t(z_t)$ in \eqref{eq:omd}-\eqref{eq:ogd}, then the approximation inequality \eqref{eq:approx-ineq} holds with $\hat{G}_F = G_F$ where $G_F$ is the uniform loss gradient ($\nabla \ell_t$) bound in Assumption~\ref{as:boundedness-grad-bregman}, i.e., 
\begin{equation*}
\|y_{t+1}-q(x_{t+1})\|_2 \leq 6 G^5 G_F^3 \eta^2\,.
\end{equation*}
\end{corollary}
\end{tcolorbox}

\begin{proof}[Proof of Proposition~\ref{prop:dist-minimizers-eta2}]

Note first that $\|\tilde{g}_t\| \leq \tilde{G}_F := G \hat{G}_F\,.$ 
Both objectives in \eqref{eq:omd}-\eqref{eq:ogd} can be written as the sum of a linear function and a strongly convex function. 
By 1-strong convexity of $R$ and uniform boundedness of $g_t$ by $\hat{G}_F$, we have for any $y \in \Y$, 
\begin{equation*}
\tilde{g}_t^{\top}(y-y_t) +\frac{1}{\eta} D_R(y\|y_t) \geq - \tilde{G}_F \|y-y_t\| + \frac{1}{2\eta} \|y-y_t\|^2\,.
\end{equation*}
Therefore, if $\|y-y_t\|_2>2 \eta \tilde{G}_F$, then the objective takes a positive value which cannot be minimal since $y=y_t$ gives zero objective value in \eqref{eq:omd}. Define the set:  
\begin{equation*}
\Y_{r,y_t}=\K \cap \{y \in \reals^d: \|y-y_t\|_2 \le r\}\,.
\end{equation*}
Then it follows that the OMD iterate in \eqref{eq:omd} satisfies: 
\begin{equation}
\label{eq:y_t+1inY_t}
y_{t+1} \in \Y_t := \Y_{2\eta \tilde{G}_F,y_t}\,. 
\end{equation}
As we want to localize both iterates \eqref{eq:omd}-\eqref{eq:ogd} in a set of size controlled by $\eta$, we now show that we also have $q(x_{t+1}) \in \Y_t$ when $y_t = q(x_t)$. Observe for this that: 
\begin{equation*}
\|q(x_{t+1}) - y_t\| = \|q(x_{t+1}) - q(x_t)\| 
\leq G \|x_{t+1} - x_t \|
\leq \eta G \|g_t\| 
\leq \eta G \hat{G}_F 
\leq 2 \eta G \hat{G}_F 
= 2 \eta \tilde{G}_F\,. 
\end{equation*}
We now write both update rules under a similar form. 
Define the function $\Phi_t: \Y \to \mathbb{R}$ by: 
\begin{equation}
\label{eq:Phi_t}
\Phi_t(y) := \tilde{g}_t^{\top}(y - y_t) + \frac{1}{2\eta} \|y-y_t\|_{\nabla^2 R(y_t)}^2\,.
\end{equation}
Note that $\Phi_t$ is $\frac{1}{\eta}$-strongly convex since $\nabla^2 \Phi_t(x) = \frac{1}{\eta} \nabla^2 R(x_t) \succeq \frac{1}{\eta} I$ as $R$ is supposed to be $1$-strongly convex.  

\begin{tcolorbox}[colframe=white!, top=2pt,left=2pt,right=2pt,bottom=2pt] 
\begin{lemma}
\label{lem:omd-rewritten}
For any $t \geq 1,$ the iterates of \eqref{eq:omd-main}-\eqref{eq:omd} can be rewritten as follows:
\begin{equation*}
y_{t+1} =\argmin_{y\in \Y_t} \left\{ \Phi_t(y) + \varepsilon_t^{\mathrm{OMD}}(y) \right\}\,, \quad \varepsilon_t^{\mathrm{OMD}}(y) := \frac{1}{\eta} \varepsilon_t(y)\,,
\end{equation*}
where $\varepsilon_t(y):= D_R(y\|y_t) - \frac12 \|y-y_t\|_{\nabla^2 R(y_t)}^2\,.$ Moreover, for $y \in \Y_t$, $|\varepsilon_t(y)| \leq \frac{4}{3} G \tilde{G}_F^3 \eta^3\,.$
\end{lemma}
\end{tcolorbox}

\begin{proof}
We have proved in \eqref{eq:y_t+1inY_t} that $y_{t+1} \in \Y_t$.  Hence, we can rewrite \eqref{eq:omd} by replacing $\Y$ by $\Y_t$: 
\begin{equation}
\label{eq:intermediate-x-rewritten}
y_{t+1} =\argmin_{y\in \Y_t} \tilde{g}_t^{\top}(y-y_t) +\frac{1}{\eta} D_R(y\|y_t). 
\end{equation}

Let $h_t :=y-y_t$. By Taylor's theorem with integral remainder, we have: 
\begin{align}
\label{eq:bregman-approx}
D_R(y\|y_t)
&=
R(y)-R(y_t)-\nabla R(y_t)^\top (y-y_t) \nonumber\\
&=
\frac12 h_t^\top \nabla^2 R(y_t)h_t
+
\varepsilon_t(y) \nonumber\\
&=
\frac12 \|h_t\|_{\nabla^2 R(y_t)}^2
+
\varepsilon_t(y),
\end{align}
where $\varepsilon_t(y)$ is the integral remainder defined as follows: 
\begin{equation*}
\varepsilon_t(y)
:= 
\frac12
\int_0^1
(1-s)^2
\nabla^3 R(y_t+s h_t)[h_t,h_t,h_t]\,ds\,,
\end{equation*}
where $\nabla^3 R$ is the third-order derivative tensor of $R$. 
Since this third-order is supposed to be uniformly bounded by $G$ over $z \in \Y$ by Assumption~\ref{as:reg-R-reparam-q-smoothness}, we have:
\begin{equation*}
|\varepsilon_t(y)|
\le
\frac{G}{6}\|h_t\|_2^3.
\end{equation*}
In particular, for \(y\in \mathcal Y_t\), i.e.
\(\|y-y_t\|_2\le 2 \eta \tilde{G}_F\), we obtain the following estimate: 
\begin{equation*}
|\varepsilon_t(y)| \leq \frac{G}{6} (2 \eta \tilde{G}_F)^3 = \frac{4}{3} G \tilde{G}_F^3 \eta^3\,.
\end{equation*}
Combining the Bregman approximation of \eqref{eq:bregman-approx} with \eqref{eq:intermediate-x-rewritten}, we obtain the desired identity using the definition of the strongly-convex function~$\Phi_t$ introduced in \eqref{eq:Phi_t}. 
\end{proof}

In the following lemma, we rewrite the OGD iterates mapped to $\Y$ via $q$ under a similar form to the update rule of OMD obtained in Lemma~\ref{lem:omd-rewritten}. 

\begin{tcolorbox}[colframe=white!, top=2pt,left=2pt,right=2pt,bottom=2pt] 
\begin{lemma}
\label{lem:ogd-rewritten}
If $(x_t)$ is the sequence of \eqref{eq:ogd-main} iterates and $y_t = q(x_t)$, then for any $t \geq 1,$
\begin{align}
\label{eq:q(u_t+1)}
z_{t+1} = q(x_{t+1}) &=\argmin_{y\in \Y_t} \left\{ \Phi_t(y) + \varepsilon_t^{\mathrm{OGD}}(y) \right\}\,,\\
\varepsilon_t^{\mathrm{OGD}}(y) &:= \tilde\varepsilon_t(y)
+ \frac{1}{\eta} \left(\frac12\|\varepsilon_t^q(y)\|_2^2 + \big\langle J_q(x_t)^{-1}(y-y_t),\,\varepsilon_t^q(y)\big\rangle
\right)\,,
\end{align}
where $\tilde\varepsilon_t(y) := g_t^\top\varepsilon_t^q(y)$ and $\varepsilon_t^q(y) := q^{-1}(y)-q^{-1}(y_t)- J_q^{-1}(x_t)(y-y_t)\,.$ Moreover, we have $\|\varepsilon_t^q(y)\|_2 \leq 2 G \tilde{G}_F^2 \eta^2$ and 
$|\tilde\varepsilon_t(y)| \leq 2 \tilde{G}_F^3 \eta^2$
\end{lemma}
\end{tcolorbox}

\begin{proof}
Recall that the \eqref{eq:ogd-main} update \eqref{eq:ogd} in the original space~$\X$ can be written as follows: 
\begin{equation}
\label{eq:update-rule-OGD}
x_{t+1}
=
\argmin_{x\in \X}
\left\{
g_t^\top(x-x_t)+\frac{1}{2\eta}\|x-x_t\|_2^2
\right\}.
\end{equation}

Recall from \eqref{eq:Phi_t} that for any $y \in \Y$: 
\begin{equation*}
\Phi_t(y) := \tilde{g}_t^{\top}(y - y_t) + \frac{1}{2\eta} \|y-y_t\|_{\nabla^2 R(y_t)}^2\,.
\end{equation*}
To prove \eqref{eq:q(u_t+1)}, we relate the linear term $g_t^\top (x-x_t)$ to $\tilde{g}_t^\top(y-y_t)$ and the quadratic term $\frac{1}{2\eta}\|x-x_t\|_2^2$ to $\frac{1}{2\eta} \|y-y_t\|_{\nabla^2 R(y_t)}^2$ respectively in \textbf{(i)} and \textbf{(ii)} below. 

\noindent\textbf{(i) Relating $g_t^\top (x-x_t)$ to $ \tilde{g}_t^\top(y-y_t).$}
Recalling that $g_t = J_q(x_t)^{\top} \tilde{g}_t$, we have:  
\begin{align}
\label{eq:inner-prod-tilde-f}
g_t^\top (x-x_t)
&= \tilde{g}_t^\top J_q(x_t)(x-x_t) \nonumber\\
&= \tilde{g}_t^\top J_q(x_t)(q^{-1}(y) - q^{-1}(y_t))\,,
\end{align}
where the last equality follows from using the identities $y = q(x)$ and $y_t = q(x_t)\,.$ 

Under Assumption~\ref{as:reg-R-reparam-q-smoothness}, applying Taylor's theorem to $q^{-1}$ around $y_t=q(x_t)$ gives, for $y=q(x)$,
\begin{equation}
\label{eq:taylor-q-1}
q^{-1}(y)-q^{-1}(y_t)
=
J_{q^{-1}}(y_t)(y-y_t)
+
\varepsilon_t^q(y)\,,
\end{equation}
where the remainder is given by: 
\begin{equation}
\label{eq:remainder-eps-t-q}
\varepsilon_t^q(y) 
= \int_0^1 (1-s)\, (y-y_t)^{\top}
\nabla^2 q^{-1}(y_t+s(y-y_t))
(y-y_t)\,ds.
\end{equation}
Since \(J_{q^{-1}}(y_t)=J_q(x_t)^{-1}\), it follows that: 
\begin{equation}
\label{eq:taylor-q-2}
q^{-1}(y)-q^{-1}(y_t) = J_q(x_t)^{-1}(y-y_t) + \varepsilon_t^q(y).
\end{equation}
Since the second derivative of \(q^{-1}\) is bounded by \(G\), we have:
\begin{equation*}
\|\varepsilon_t^q(y)\|_2
\le
\frac{G}{2}\|y-y_t\|_2^2.
\end{equation*}
In particular, for $y\in \Y_t = \mathcal Y_{2\eta \tilde{G}_F,y_t}$, i.e. $y \in \Y$ s.t.
$\|y-y_t\|_2\le 2\eta \tilde{G}_F$, we obtain: 
\begin{equation}
\label{eq:estimate-vareps-tqy}
\|\varepsilon_t^q(y)\|_2 \leq \frac{G}{2} (2\eta \tilde{G}_F)^2 = 2 G \tilde{G}_F^2 \eta^2\,.
\end{equation}

Using \eqref{eq:taylor-q-2} in \eqref{eq:inner-prod-tilde-f} yields:  
\begin{align}
\label{eq:term-inner-prod-rewriting}
g_t^\top (x-x_t)
&= \tilde{g}_t^\top J_q(x_t)(q^{-1}(y) - q^{-1}(y_t)) \nonumber\\
&= \tilde{g}_t^\top(y-y_t) + \tilde{g}_t^\top J_q(x_t)\,\varepsilon_t^q(y) \nonumber\\
&= \tilde{g}_t^\top(y-y_t) + \tilde\varepsilon_t(y),
\end{align}
where 
$\tilde\varepsilon_t(y) :=
\tilde{g}_t^\top J_q(x_t)\,\varepsilon_t^q(y)
= g_t^{\top} \varepsilon_t^q(y)$. Therefore, using \eqref{eq:estimate-vareps-tqy}, we have for any $y \in \Y_t,$ 
\begin{equation*}
|\tilde\varepsilon_t(y)| \leq \|g_t\| \cdot \|\varepsilon_t^q(y)\| \leq G (2\hat{G}_F \tilde{G}_F^2 \eta^2) = 2 \tilde{G}_F^3 \eta^2\,,
\end{equation*}
where the last equality stems from recalling that $\tilde{G}_F = G \hat{G}_F\,.$ 

\noindent\textbf{(ii) Relating $\frac{1}{2\eta}\|x-x_t\|_2^2$ to $\frac{1}{2\eta} \|y-y_t\|_{\nabla^2 R(y_t)}^2.$} For any $x \in \X$, $y = q(x),$ given that $y_t = q(x_t)$,  
\begin{align}
\label{eq:u-u_t-square}
\frac{1}{2\eta}\|x-x_t\|_2^2
&= \frac{1}{2\eta}\|q^{-1}(y)-q^{-1}(y_t)\|_2^2 \nonumber\\
&= \frac{1}{2\eta}\|J_q(x_t)^{-1}(y-y_t) + \varepsilon_t^q(y)\|_2^2 \nonumber\\
&= \frac{1}{2\eta}\|J_q(x_t)^{-1}(y-y_t)\|_2^2
+ \frac{1}{\eta} \left(\frac12\|\varepsilon_t^q(y)\|_2^2
+\big\langle J_q(x_t)^{-1}(y-y_t),\,\varepsilon_t^q(y)\big\rangle\right).
\end{align}
Using Assumption~\ref{as:reparam} relating $J_q$ and $\nabla^2R$ and the initial coupling $y_t = z_t = q(x_t)$, we have: 
\begin{equation*}
\big[\nabla^2R(y_t)\big]^{-1}=J_q(x_t)J_q(x_t)^\top,
\qquad
\nabla^2R(y_t)=[J_q(x_t)^{\top}]^{-1}J_q(x_t)^{-1}\,.
\end{equation*}
Therefore we can write: 
\begin{align}
\label{eq:Jq-1(x-x_t)}
\|J_q(x_t)^{-1}(y-y_t)\|_2^2
&= (y-y_t)^{\top} [J_q(x_t)^{-1}]^{\top} J_q(x_t)^{-1}(y-y_t) \nonumber\\
&= (y-y_t)^{\top} \nabla^2R(y_t)(y-y_t) \nonumber\\
&= \|y-y_t\|^2_{\nabla^2R(y_t)}.
\end{align}
Combining \eqref{eq:u-u_t-square} and \eqref{eq:Jq-1(x-x_t)} yields: 
\begin{equation*}
\label{eq:term-reg-rewriting}
\frac{1}{2\eta}\|x-x_t\|_2^2 = \frac{1}{2\eta} \|y-y_t\|^2_{\nabla^2R(y_t)} 
+ \frac{1}{\eta} \left(\frac12\|\varepsilon_t^q(y)\|_2^2
+\big\langle J_q(x_t)^{-1}(y-y_t),\,\varepsilon_t^q(y)\big\rangle\right)\,.
\end{equation*}

Given the OGD update rule \eqref{eq:update-rule-OGD}, the change of variable $y = q(x)$ and \eqref{eq:term-inner-prod-rewriting}, \eqref{eq:term-reg-rewriting} proven in \textbf{(i)} and \textbf{(ii)} respectively, we have shown that 
\begin{equation*}
q(x_{t+1})
=
\argmin_{y\in \Y_t}
\left\{
\Phi_t(y)+\varepsilon_t^{\mathrm{OGD}}(y)
\right\},
\end{equation*}
where the error term $\varepsilon_t^{\mathrm{OGD}}(y)$ is defined as follows: 
\begin{equation}
\label{eq:eps-OGD}
\varepsilon_t^{\mathrm{OGD}}(y)
:=
\tilde\varepsilon_t(y)
+
\frac{1}{\eta}
\left(
\frac12\|\varepsilon_t^q(y)\|_2^2
+
\big\langle J_q(x_t)^{-1}(y-y_t),\,\varepsilon_t^q(y)\big\rangle
\right).
\end{equation}
\end{proof}

So far, we have shown in Lemmas~\ref{lem:omd-rewritten} and~\ref{lem:ogd-rewritten} respectively the following similar update rule forms: 
\begin{align}
y_{t+1}
&=
\argmin_{y\in \Y_t}\{\Phi_t(y)+\varepsilon_t^{\mathrm{OMD}}(y)\},
\label{eq:min1}
\\
z_{t+1} = q(x_{t+1})
&=
\argmin_{y\in \Y_t}\{\Phi_t(y)+\varepsilon_t^{\mathrm{OGD}}(y)\}.
\label{eq:min2}
\end{align}
In the next lemma, we use first-order conditions together with strong convexity of the function~$\Phi_t$ to show that the distance between the minimizers can be upperbounded by $\eta$ times the distance between the gradients of the errors in each one of the updates (for OMD and OGD). Note that our proof technique here is different from the strategy adopted in \citet{ghai-lu-hazan22} (Lemma 4) as we do not control the distance between objectives and translate it to the distance between minimizers using strong convexity. 

\begin{tcolorbox}[colframe=white!, top=2pt,left=2pt,right=2pt,bottom=2pt] 
\begin{lemma} 
\label{lem:dist-minimizers-eta-times-norm}
For any $t \geq 1$, if $y_t = q(x_t)$, then 
\begin{equation*}
\|y_{t+1} - q(x_{t+1})\| \le \eta\,
\left\|\nabla\varepsilon_t^{\mathrm{OGD}}(q(x_{t+1})) - \nabla\varepsilon_t^{\mathrm{OMD}}(y_{t+1})\right\|\,.
\end{equation*}
\end{lemma}
\end{tcolorbox}

To prove this result, we will need the following technical lemma stating monotonicity of normal cones of nonempty closed convex sets, which is a standard result in convex analysis. 

\begin{lemma}[Monotonicity of the normal cone]
\label{lem:normal-cone-monotonicity}
Let \(K\subset \mathbb R^d\) be a nonempty closed convex set. For any
\(x,y\in K\), let
\[
N_K(x):=\{v\in\mathbb R^d:\langle v,z-x\rangle\le 0
\,, \quad \forall z\in K\}
\]
denote the normal cone of \(K\) at \(x\). Then, for any
\(v_x\in N_K(x)\) and \(v_y\in N_K(y)\),
\[
\langle v_x-v_y,x-y\rangle \ge 0.
\]
\end{lemma}

\begin{proof}
Since \(v_x\in N_K(x)\) and \(y\in K\), the definition of the normal cone gives $\langle v_x,y-x\rangle\le 0,$
or equivalently, $\langle v_x,x-y\rangle\ge 0.$
Similarly, since \(v_y\in N_K(y)\) and \(x\in K\), we have 
$\langle v_y,x-y\rangle\le 0.$
Therefore, it follows from the above inequalities that: 
\begin{equation*}
\langle v_x-v_y,x-y\rangle
=
\langle v_x,x-y\rangle
-
\langle v_y,x-y\rangle
\ge 0\,.
\end{equation*}
\end{proof}

\begin{proof}[Proof of Lemma~\ref{lem:dist-minimizers-eta-times-norm}]
Let $\delta_{t+1}:=y_{t+1}-q(x_{t+1})$ as a shorthand notation.  
From first-order optimality conditions for \eqref{eq:min1} and \eqref{eq:min2},
there exist $v_{t+1}^{\mathrm{OMD}} \in N_{\Y_t}(y_{t+1}), v_{t+1}^{\mathrm{OGD}} \in N_{\Y_t}(q(x_{t+1}))$ s.t. 
\begin{align*}
\nabla \Phi_t(y_{t+1})+\nabla \varepsilon_t^{\mathrm{OMD}}(y_{t+1})+v_{t+1}^{\mathrm{OMD}} &= 0\,,\\
\nabla \Phi_t(q(x_{t+1}))+\nabla \varepsilon_t^{\mathrm{OGD}}(q(x_{t+1}))+v_{t+1}^{\mathrm{OGD}} &= 0\,.
\end{align*}
Subtracting both identities and taking inner product with $\delta_{t+1}$ gives
\begin{multline}
\big\langle \nabla\Phi_t(y_{t+1})-\nabla\Phi_t(q(x_{t+1})),\,\delta_{t+1}\big\rangle
+ \big\langle \nabla\varepsilon_t^{\mathrm{OMD}}(y_{t+1})
- \nabla\varepsilon_t^{\mathrm{OGD}}(q(x_{t+1})),\,\delta_{t+1}\big\rangle \\
+ \big\langle v_{t+1}^{\mathrm{OMD}}-v_{t+1}^{\mathrm{OGD}}, \,\delta_{t+1} \big\rangle
=0.
\label{eq:master}
\end{multline}
By $\frac{1}{\eta}$-strong convexity of $\Phi_t$, we have: 
\begin{equation}
\label{eq:sc-phi_t-ineq}
\big\langle \nabla\Phi_t(y_{t+1})-\nabla\Phi_t(q(x_{t+1})), \,\delta_{t+1}\big\rangle
\ge \frac1\eta \|\delta_{t+1}\|^2.
\end{equation}
By monotonicity of normal cones of convex sets (Lemma~\ref{lem:normal-cone-monotonicity} above),
we can also control the term involving the normal cone directions as follows: 
\begin{equation}
\label{eq:normal-cone-monotonicity-ineq}
\big\langle v_{t+1}^{\mathrm{OMD}}-v_{t+1}^{\mathrm{OGD}}, \,\delta_{t+1} \big\rangle \ge 0.
\end{equation}
Plugging \eqref{eq:sc-phi_t-ineq} and \eqref{eq:normal-cone-monotonicity-ineq} into \eqref{eq:master}, we obtain: 
\begin{equation*}
\frac1\eta\|\delta_{t+1}\|^2
\le 
\left\| \nabla\varepsilon_t^{\mathrm{OGD}}(q(x_{t+1})) - \nabla\varepsilon_t^{\mathrm{OMD}}(y_{t+1}) \right\|
\cdot \|\delta_{t+1}\|.
\end{equation*}
Therefore, it follows that the distance~$\delta_{t+1}$ between the two minimizers is upperbounded as follows: 
\[
\|\delta_{t+1}\|
\le
\eta\,
\left\|
\nabla\varepsilon_t^{\mathrm{OGD}}(q(x_{t+1}))
-
\nabla\varepsilon_t^{\mathrm{OMD}}(y_{t+1})
\right\|.
\]
\end{proof}

It remains to bound the last norm in the right-hand side of Lemma~\ref{lem:dist-minimizers-eta-times-norm} and show it is also of order $\eta$. We show that each one of the norms of the gradients $\|\nabla\varepsilon_t^{\mathrm{OMD}}(y_{t+1})\|$ and $\|\nabla\varepsilon_t^{\mathrm{OGD}}(q(x_{t+1}))\|$ are of order $\eta$ in the two following lemmas. 

\begin{tcolorbox}[colframe=white!, top=2pt,left=2pt,right=2pt,bottom=2pt] 
\begin{lemma}
\label{lem:norm-grad-eps-OMD}
$\|\nabla \varepsilon_t^{\mathrm{OMD}}(y_{t+1})\| \leq 2 G \tilde{G}_F^2 \eta\,.$
\end{lemma}
\end{tcolorbox}

\begin{proof}
Recall that $\varepsilon_t^{\mathrm{OMD}}(y)=\frac1\eta \varepsilon_t(y),$
where: 
\begin{equation*}
\varepsilon_t(y)
=
D_R(y\|y_t)-\frac12(y-y_t)^\top \nabla^2R(y_t)(y-y_t)\,.
\end{equation*}
Differentiating the above error function w.r.t. $y$ yields: 
\begin{equation*}
\nabla \varepsilon_t(y)
=
\nabla R(y)-\nabla R(y_t)-\nabla^2R(y_t)(y-y_t).
\end{equation*}
Using third-order smoothness of $R$ (bounded third derivative) as supposed in Assumption~\ref{as:reg-R-reparam-q-smoothness},
\begin{equation}
\label{eq:bound-grad-eps-t}
\|\nabla \varepsilon_t(y_{t+1})\|
\le
\frac{G}{2}\|y_{t+1}-y_t\|^2.
\end{equation}
Since $y_{t+1}\in \Y_t$, we have $\|y_{t+1}-y_t\| \leq 2\eta \tilde{G}_F$ and it follows from \eqref{eq:bound-grad-eps-t} that: 
\begin{equation*}
\|\nabla \varepsilon_t(y_{t+1})\| \leq \frac{G}{2} (2\eta \tilde{G}_F)^2 = 2 G \tilde{G}_F^2 \eta^2\,.
\end{equation*}
Since $\varepsilon_t^{\mathrm{OMD}}(y)=\frac1\eta \varepsilon_t(y),$ we conclude that: 
\begin{equation*}
\|\nabla \varepsilon_t^{\mathrm{OMD}}(y_{t+1})\|
= \frac 1 \eta \|\nabla \varepsilon_t(y_{t+1})\| 
\leq 2 G \tilde{G}_F^2 \eta\,.
\end{equation*}
\end{proof}

\begin{tcolorbox}[colframe=white!, top=2pt,left=2pt,right=2pt,bottom=2pt] 
\begin{lemma}
\label{lem:norm-grad-eps-OGD}
$\|\nabla\varepsilon_t^{\mathrm{OGD}}(q(x_{t+1}))\| \leq G^5 (3\hat{G}_F^2 + \hat{G}_F^3 \eta) \eta.$
\end{lemma}
\end{tcolorbox}

\begin{proof}
Recall from \eqref{eq:eps-OGD} that: 
\[
\varepsilon_t^{\mathrm{OGD}}(y)
=
\tilde\varepsilon_t(y)
+
\frac1\eta
\left(
\frac12\|\varepsilon_t^q(y)\|_2^2
+
\big\langle J_q(x_t)^{-1}(y-y_t),\,\varepsilon_t^q(y)\big\rangle
\right).
\]
Differentiating this error term w.r.t. $y$ gives: 
\begin{equation}
\label{eq:grad-eps-OGD}
\nabla\varepsilon_t^{\mathrm{OGD}}(y)
=
\nabla \tilde\varepsilon_t(y)
+
\frac1\eta J_{\varepsilon_t^q}(y)^\top \varepsilon_t^q(y)
+
\frac1\eta\Big(
[J_q(x_t)^{-1}]^\top \varepsilon_t^q(y)
+
J_{\varepsilon_t^q}(y)^\top J_q(x_t)^{-1}(y-y_t)
\Big).
\end{equation}
Since $\tilde\varepsilon_t(y) := g_t^\top \varepsilon_t^q(y)$, we have: 
\begin{equation*}
\nabla \tilde\varepsilon_t(y)
=
g_t^\top J_{\varepsilon_t^q}(y)\,, \quad \|\nabla \tilde\varepsilon_t(y)\| \leq \hat{G}_F \|J_{\varepsilon_t^q}(y)\|\,.
\end{equation*}

We now prove the following estimates: 
\begin{enumerate}[label=(\roman*)]
\item \label{eq:i} $\|q(x_{t+1})-y_t\| \leq  G \hat{G}_F \eta\,,$
\item \label{eq:ii} $\|J_{\varepsilon_t^q}(q(x_{t+1}))\| \leq  G^2 \hat{G}_F \eta,$
\item \label{eq:iii} $\|\varepsilon_t^q(q(x_{t+1}))\| \leq \frac{G^3 \hat{G}_F^2}{2} \eta^2\,.$ 
\end{enumerate}

\noindent\textbf{Proof of \ref{eq:i}.}  We control this first error term as follows: 
\begin{align}
\|q(x_{t+1})-y_t\|
&= \|q(x_{t+1})-q(x_t)\|
&&\hfill \text{(since $y_t = q(x_t)$)}  \nonumber\\
&\leq G \|x_{t+1} - x_t\|
&&\hfill \text{($q$ is $G$-Lipschitz by Assumption~\ref{as:reg-R-reparam-q-smoothness})} \nonumber\\
&= G\|\Pi_{\X_{\delta}}(x_t-\eta g_t) - x_t\|
&&\hfill \text{(by definition of $x_{t+1}$ in \eqref{eq:ogd-main} or BOGD)} \nonumber\\
&= G \|\Pi_{\X_{\delta}}(x_t-\eta g_t) - \Pi_{\X_{\delta}}(x_t)\|
&&\hfill \text{(since $x_t\in\X_{\delta}$, $\Pi_{\X_{\delta}}(x_t)=x_t$)} \nonumber\\
&\leq \eta G \|g_t\| 
&&\hfill \text{(by nonexpansiveness of the projection)} \nonumber\\
&\leq \eta G \hat{G}_F = \eta \tilde{G}_F \,.
&&\hfill \text{(using Assumption~\ref{as:boundedness-grad-bregman})}
\end{align}

\noindent\textbf{Proof of \ref{eq:ii}.} Recall from \eqref{eq:taylor-q-1} that for any $y=q(x), x \in \X\,,$ 
\begin{equation}
\label{eq:eps-t-q-taylor}
\varepsilon_t^q(y) = q^{-1}(y)-q^{-1}(y_t) - J_{q^{-1}}(y_t)(y-y_t)\,.
\end{equation}
Differentiating w.r.t. $y$ yields: 
\begin{equation*}
J_{\varepsilon_t^q}(q(x_{t+1})) = J_{q^{-1}}(q(x_{t+1})) - J_{q^{-1}}(y_t)\,.
\end{equation*}
Using Assumption~\ref{as:reg-R-reparam-q-smoothness}, we obtain: 
\begin{equation*}
\|J_{\varepsilon_t^q}(q(x_{t+1}))\| = \|J_{q^{-1}}(q(x_{t+1})) - J_{q^{-1}}(y_t)\| \leq G \|q(x_{t+1}) - y_t\| \leq \eta G^2 \hat{G}_F\,,
\end{equation*}
where the last step uses the first estimate \ref{eq:i} proved above.

\noindent\textbf{Proof of \ref{eq:iii}.} Recalling \eqref{eq:eps-t-q-taylor} and using boundedness of the second-order derivative of $q^{-1}$ (Assumption~\ref{as:reg-R-reparam-q-smoothness}), we immediately obtain: 
\begin{equation*}
\|\varepsilon_t^q(q(x_{t+1}))\| \leq \frac{G}{2} \cdot \|q(x_{t+1}) - y_t\|^2 \leq \frac{G^3 \hat{G}_F^2}{2} \eta^2\,,
\end{equation*}
where the last estimate follows from using the first proven estimate \ref{eq:i}: $\|q(x_{t+1})-y_t\|\leq  G \hat{G}_F \eta\,.$

We conclude the proof of Lemma~\ref{lem:norm-grad-eps-OGD} by using all the estimates \ref{eq:i} to \ref{eq:iii} to bound the norm of the gradient~$\nabla\varepsilon_t^{\mathrm{OGD}}(x)$ in \eqref{eq:grad-eps-OGD} as follows:
\begin{align}
\label{eq:grad-eps-ogd}
\|\nabla\varepsilon_t^{\mathrm{OGD}}(q(x_{t+1}))\| 
&\leq \hat{G}_F \|J_{\varepsilon_t^q}(q(x_{t+1}))\| 
+ \frac{1}{\eta} \|J_{\varepsilon_t^q}(q(x_{t+1}))\| \cdot \|\varepsilon_t^q(q(x_{t+1}))\| \nonumber\\
&+ \frac{1}{\eta} G \|\varepsilon_t^q(q(x_{t+1}))\| + \frac{1}{\eta} G \|J_{\varepsilon_t^q}(q(x_{t+1}))\| \cdot \|q(x_{t+1})-y_t\| \nonumber\\
&\leq G^2  \hat{G}_F^2 \eta + \frac 12 G^5 \hat{G}_F^3 \eta^2 + \frac 12 G^4 \hat{G}_F^2 \eta + G^4 \hat{G}_F^2 \eta\,\nonumber\\
&\leq G^5 (3\hat{G}_F^2 + \hat{G}_F^3 \eta) \eta\,,
\end{align}
where the last inequality follows from using the assumption $G \geq 1$ which is without loss of generality (otherwise replace the constant by 1). 
\end{proof}

\noindent\textbf{End of Proof of Proposition~\ref{prop:dist-minimizers-eta2}.} 
Using Lemma~\ref{lem:dist-minimizers-eta-times-norm}, we have: 
\begin{equation*}
\|y_{t+1} - q(x_{t+1})\| \le \eta\,
\left\|\nabla\varepsilon_t^{\mathrm{OGD}}(q(x_{t+1})) - \nabla\varepsilon_t^{\mathrm{OMD}}(y_{t+1})\right\|\,.
\end{equation*}
Then it follows from using Lemma~\ref{lem:norm-grad-eps-OMD} and Lemma~\ref{lem:norm-grad-eps-OGD} that: 
\begin{align}
\|y_{t+1}-q(x_{t+1})\| = 
\|\delta_{t+1}\|
&\le
\eta\cdot
\left(
\|\nabla\varepsilon_t^{\mathrm{OGD}}(q(x_{t+1}))\|
+
\|\nabla\varepsilon_t^{\mathrm{OMD}}(y_{t+1})\|
\right)\nonumber\\
&\leq (G^5 (3\hat{G}_F^2 + \hat{G}_F^3 \eta) + 2 G \tilde{G}_F^2) \eta^2 \nonumber\\
&= (G^5 (3\hat{G}_F^2 + \hat{G}_F^3 \eta) + 2 G^3 \hat{G}_F^2)\eta^2\nonumber\\
&\leq G^5 (5\hat{G}_F^2 + \hat{G}_F^3 \eta) \eta^2\,,
\end{align}
which concludes the proof. 
\end{proof}

\subsection{Proof of Theorem~\ref{thm:main} using Lemma~\ref{prop:dist-minimizers-eta2-main}}

The following lemma controls the regret of an approximate OMD algorithm. 
\begin{lemma}[Lemma 6 in \citet{ghai-lu-hazan22}]
\label{lem2} 
Suppose Assumptions~\ref{as:reg-R-reparam-q-smoothness} and ~\ref{as:boundedness-grad-bregman} hold and let Algorithm $\A$ update the sequence $(z_t)$ as follows:
$$
    z_{t+1}=r_{t+1} +\argmin_{y\in \Y} \left\{ \nabla h_t(y_t)^{\top}(y-y_t) +\frac{1}{\eta} D_R(y\|y_t) \right\}\,, 
$$
where $\|r_{t+1}\|_2 \le C_{\eta}$. Then the regret of Algorithm $\A$ is upper bounded as follows: 
$$
R_T(\A) \le \frac{C_{\eta} T G}{\eta} + \frac{D_1}{\eta} +\frac{\eta G^2 T}{2}\,.
$$
\end{lemma}


\noindent\textbf{End of Proof of Theorem~\ref{thm:main}.} 
The rest of the proof follows the same lines as the proof of \citet{ghai-lu-hazan22} using our tighter estimate obtained in Corollary~\ref{cor:approx-lemma-det}. We provide a full proof for completeness. 
\begin{proof}
The \eqref{eq:ogd-main} iterates mapped to the hidden space $\Y$ can be seen as a perturbed version of OMD since: 
\begin{equation*}
z_{t+1} = q(x_{t+1}) = q(x_{t+1}) - y_{t+1} + y_{t+1} = r_{t+1} + \argmin_{y\in \Y} \nabla h_t(y_t)^{\top}(y-y_t) +\frac{1}{\eta} D_R(y\|y_t)\,, 
\end{equation*}
where $r_{t+1} := q(x_{t+1}) - y_{t+1}\,.$ It follows from Corollary~\ref{cor:approx-lemma-det} that $\|r_{t+1}\| \leq C_{\eta} := 6 G^5 G_F^3 \eta^2\,.$ Using now Lemma~\ref{lem2}, we obtain the following regret bound for the \eqref{eq:ogd-main} iterates $(z_t)$ in the hidden space $\Y$:
\begin{equation*}
R_T = \sum_{t=1}^T h_t(z_t) - \min_{z\in\Y}\sum_{t=1}^T h_t(z) \le \frac{C_{\eta} T G}{\eta} + \frac{D_1}{\eta} +\frac{\eta G^2 T}{2} 
\leq \frac{D_1}{\eta} + 7 G^6 G_F^3  \eta T\,.
\end{equation*}
Optimizing the stepsize $\eta$ to minimize the upperbound yields: 
\begin{equation*}
R_T \leq \sqrt{7 D_1 G^6 G_F^3 T}\,, 
\end{equation*}
with stepsize $\eta = \sqrt{\frac{D_1}{7 G^6 G_F^3 T}}\,.$ 
\end{proof}

\section{Proofs for Section~\ref{sec:discussion-hessian-comp-as}}

\subsection{Proof of Proposition~\ref{prop:matrix-field-eq-Hessian}}

The proof relies on the following classical result from vector calculus. 
\begin{tcolorbox}[colframe=white!, top=2pt,left=2pt,right=2pt,bottom=2pt] 
\begin{proposition}
\label{prop:existence-grad-potential}
Let $\mathcal{U} \subseteq \mathbb{R}^d$ be convex\footnote{A simple connected domain suffices here. Convexity is enough for our purpose.} and let $F: \mathcal{U} \to \mathbb{R}^d $ be a $C^1$ vector field. Denoting $F = (F_1, \cdots, F_d)$ using its coordinate functions, if $\partial_{x_j} F_k(x) = \partial_{x_k} F_j(x)$ for all $j,k \in [d]$ and all $x \in \mathcal{U}$, then there exists a scalar function $\phi: \mathbb{R}^d \to \mathbb{R}$ such that $F = \nabla \phi\,.$ 
\end{proposition}
\end{tcolorbox}

\begin{proof}[Proof of Prop.~\ref{prop:matrix-field-eq-Hessian}]
For each $i \in [d]$, define the vector field $M^{(i)}$ of the $i$th row of the matrix field $M$: 
\begin{equation*}
M^{(i)}(x) := (M_{i,1}(x), \cdots, M_{i,d}(x))\,.
\end{equation*}
The cross-derivatives conditions $\partial_{x_k} M_{ij}(x) = \partial_{x_j} M_{ik}(x)\,,\forall i, j, k \in [d], \forall x \in K\,,$ guarantee that $M^{(i)}$ has also matching cross-partial derivatives in the sense of Proposition~\ref{prop:existence-grad-potential}. Therefore, it follows from Proposition~\ref{prop:existence-grad-potential} that there exists a scalar function $g_i: \mathbb{R}^d \to \mathbb{R}$ such that $\nabla g_i = F^{(i)}$ for all $i \in [d]\,.$ 

We obtain a vector field $g = (g_1, \cdots, g_d)$ whose Jacobian $J_g$ coincides with $M$, i.e. $J_g = M$ since $M^{(i)}$ are the rows of the matrix field $M$. Then, observe that $M$ is symmetric because $J_q J_q^{\top}$ is symmetric and its inverse is also symmetric. Hence, the Jacobian $J_g = M$ is also symmetric and it follows that $g$ itself has matching cross partial derivatives. Applying Proposition~\ref{prop:existence-grad-potential} again to $g$ yields the existence of a scalar function $R$ such that $g = \nabla R\,.$ We have proved that $M = J_g$ and $g = \nabla R$ which implies that $M = \nabla^2 R$, concluding the proof. 
\end{proof}

\subsection{Detailed derivations for examples}
\label{appx:detailed-examples}

\paragraph{Example 1: Affine mixing of a separable reparameterization.}
Let
\[
    q(x)=A\,s(x)+b,
    \qquad
    s(x)=\bigl(s_1(x_1),\dots,s_d(x_d)\bigr),
\]
where \(A\in\mathbb R^{d\times d}\) is invertible, \(b\in\mathbb R^d\), and each
\(s_i:\mathbb R\to\mathbb R\) is a \(C^2\) diffeomorphism. 
Writing
\[
    D(x):=\operatorname{diag}\bigl(s_1'(x_1),\dots,s_d'(x_d)\bigr),
\]
it follows that the Jacobian of $q$ at $x$ and its outer product can be written as: 
\[
    J_q(x)=A D(x),
    \qquad
    J_q(x)J_q(x)^\top=A D(x)^2 A^\top.
\]
Therefore, for \(z=q(x)\),
\[
    M(z)
    =
    \left[J_q(q^{-1}(z))J_q(q^{-1}(z))^\top\right]^{-1}
    =
    A^{-T}D(x)^{-2}A^{-1},
    \qquad x=q^{-1}(z).
\]

Let $\xi:=A^{-1}(z-b).$ 
Since $z=A s(x)+b$, we have $\xi=s(x)$, hence \(x_i=s_i^{-1}(\xi_i)\). Define
\[
    m_i(\xi_i):=
    \frac{1}{\bigl(s_i'(s_i^{-1}(\xi_i))\bigr)^2}.
\]
Using this notation, we can rewrite $M(z)$ as follows: 
\[
    M(z)
    =
    A^{-T}\operatorname{diag}\bigl(m_1(\xi_1),\dots,m_d(\xi_d)\bigr)A^{-1},
    \qquad
    \xi=A^{-1}(z-b).
\]

This matrix field is explicitly Hessian. Indeed, define \(r_i\) by
\[
    r_i''(t)=m_i(t)
    =
    \frac{1}{\bigl(s_i'(s_i^{-1}(t))\bigr)^2},
\]
and set
\[
    R(z):=\sum_{i=1}^d r_i\bigl((A^{-1}(z-b))_i\bigr).
\]
Since \(A^{-1}(z-b)\) is affine in \(z\),
\[
    \nabla^2 R(z)
    =
    A^{-T}\operatorname{diag}\bigl(r_1''(\xi_1),\dots,r_d''(\xi_d)\bigr)A^{-1}
    =
    M(z).
\]
Thus $M$ satisfies the Hessian-compatibility condition. 
In particular, the compatibility holds although $J_q(x)=AD(x)$ is not diagonal unless $A$ is diagonal.

\paragraph{Example 2: Rank-one nonlinear mixing.}
Fix \(a\in\mathbb R^d\setminus\{0\}\) and let
\begin{equation*}
    q(x)=x+h(a^\top x)a,
\end{equation*}
where \(h:\mathbb R\to\mathbb R\) is \(C^2\). 
Writing \(r:=a^\top x\) and \(s:=\|a\|^2\), we have
\[
    J_q(x)=I+h'(r)aa^\top.
\]
This Jacobian is generally non-diagonal whenever at least two components of \(a\) are nonzero.

The matrix \(J_q(x)\) acts as multiplication by \(1+s h'(r)\) in the direction \(a\), and as the identity on \(a^\perp\). 
Assume \(1+s h'(r)\neq 0\) on the domain, so that \(q\) is locally invertible. 
Then \(J_q(x)J_q(x)^\top\) has eigenvalue \((1+s h'(r))^2\) along \(a\) and eigenvalue \(1\) on \(a^\perp\). 
Consequently, the inverse metric field has the form
\[
    M(z)=I+\beta(a^\top z)\,aa^\top,
    \qquad z=q(x),
\]
for a scalar function \(\beta\). To identify it, note that
\[
    a^\top z
    =
    a^\top q(x)
    =
    r+s h(r).
\]
Let \(g(r):=r+s h(r)\), and assume \(g\) is invertible on the relevant interval. Then
\[
    r=g^{-1}(a^\top z),
\]
and the eigenvalue of \(M(z)\) along \(a\) gives
\[
    1+s\,\beta(a^\top z)
    =
    \bigl(1+s h'(r)\bigr)^{-2}.
\]
Thus we can deduce an expression for $\beta$ as follows: 
\[
    \beta(t)
    =
    \frac{\bigl(1+s h'(g^{-1}(t))\bigr)^{-2}-1}{s}.
\]

Again \(M\) is explicitly Hessian. Let \(\psi:\mathbb R\to\mathbb R\) satisfy $\psi''(t)=\beta(t)$, 
and define: 
\[
    R(z):=\frac12\|z\|^2+\psi(a^\top z).
\]
Then, it follows from the above definition that the Hessian can be written as follows: 
\[
    \nabla^2R(z)
    =
    I+\psi''(a^\top z)aa^\top
    =
    I+\beta(a^\top z)aa^\top
    =
    M(z).
\]
Hence the Hessian-compatibility condition holds, despite the generally non-diagonal Jacobian \(J_q(x)\).

\subsection{Proof of Theorem~\ref{thm:linear-reg-lower-bound}: Linear Regret Lower Bound for OGD}
\label{app:proof-thm-linear-reg-lower-bound}

\subsection{Detailed proof sketch of Theorem~\ref{thm:linear-reg-lower-bound}} 
\label{app:detailed-proof-sketch-thm-linear-reg-lower-bound}

The proof consists of several steps. We give an overview of each one of them below. 
The complete proof with all the details can be found in Section~\ref{app:proof-thm-linear-reg-lower-bound}.\\ 

\noindent\textbf{Step 1. Parameterization $q$ and induced metric}. 
Let $\mathcal{X} := [0,\log 3] \times [\frac \pi 6, \frac \pi 3] \subset \mathbb{R}^2.$ Define the parameterization function $q: \mathcal{X} \to \mathbb{R}^2$ for any $x = (x_1, x_2) \in \mathcal{X}$ as follows: 
$q(x_1,x_2) := e^{x_1} (\cos x_2, \sin x_2)\,.$ 
One can then easily verify that $J_q(x) J_q(x)^{\top} = e^{2x_1} I_2$ where $I_2$ is the $2\times 2$ identity matrix. 
Moreover, since $x_1 \geq 0$, $e^{2x_1} \geq 1$ and hence $J_q(x) J_q(x)^{\top} \succeq I_2$. 
$z := q(x)$, we have $\|z\| = \|q(x)\| = e^{x_1}$ and therefore $G(z) := M(z)^{-1} = \frac{1}{\|z\|^2} I_2\,.$ 

\noindent\textbf{Step 2: OGD dynamics in the hidden space.} To construct our adversarial sequence of losses, we first write the OGD dynamics in the hidden space $\Y := q(\mathcal{X}).$ 
We will consider linear functions $h_t: \Y \to \mathbb{R}$ defined for every $z \in \mathcal{Z}$ 
by $h_t(z) = \langle s_t, z \rangle$ where $(s_t)$ is an adversarial sequence that will be specified later in an adversarial way. 
It follows that for any $x \in \mathcal{X}$, $\ell_t(x) = h_t(q(x)) = \langle s_t, q(x) \rangle\,,$ and $\nabla \ell_t(x) = J_q(x)^{\top} s_t$.
The OGD update rule on the sequence $(f_t)$ can then be written as follows: 
\begin{equation}
\label{eq:ogd-update-x-main}
x_{t+1} = x_t - \eta \nabla \ell_t(x_t) = x_t - \eta J_q(x_t)^{\top} s_t\,.
\end{equation}
Now defining the sequence $(z_t)$ by $z_t := q(x_t)$ for any $t$ in the hidden space $\mathcal{Z}$, we obtain the following using Taylor's theorem with exact remainder: 
\begin{equation}
\label{eq:1-taylor-main}
z_{t+1} = q(x_{t+1}) = q(x_t) + J_q(x_t) (x_{t+1} - x_t) + r_t = z_t - \eta G(z_t)^{-1} s_t + r_t\,,
\end{equation}
where $\|r_t\| = \mathcal{O }( \eta^2)\,.$ 
To simplify the OGD dynamics in the hidden space, we choose $s_t = - G(z_t) v_t$ where $(v_t)$ is an adversarial sequence that will be chosen later. 
Given this choice, we obtain the following OGD dynamics in the hidden space: 
\begin{equation}
\label{eq:ogd-hidden-space-main}
z_{t+1} = z_t + \eta v_t + r_t\,.
\end{equation}

\noindent\textbf{Step 3. Regret as an approximate discrete path sum.} 
For a fixed comparator $u \in \mathcal{X}$, let $z_u := q(u).$ Then we have by setting $s_t = - G(z_t) v_t$: 
\begin{equation*}
\ell_t(x_t) - \ell_t(u) 
= h_t(z_t) - h_t(z_u)
= \langle s_t, z_t - z_u \rangle 
=  - \langle G(z_t) v_t, z_t - z_u \rangle 
= - \langle G(z_t) (z_t - z_u), v_t  \rangle\,,
\end{equation*}
where the last step uses symmetry of $G(z_t)$. 
Define the vector field $F_u: \mathcal{Z} \to \mathbb{R}^2$ for any $z \in \mathcal{Z}$ by:
\begin{equation}
\label{eq:def-F_u-main}
F_u(z) := G(z) (z - z_u)\,.
\end{equation}
Using the above with the hidden-space dynamics \eqref{eq:ogd-hidden-space-main} 
and the estimate $\|r_t\| = \mathcal{O}(\eta^2)$, we obtain:  
\begin{align}
    \ell_t(x_t) - \ell_t(u) = - \frac{1}{\eta} \langle F_u(z_t), z_{t+1} - z_t\rangle +  \mathcal{O}(\eta)\,.  
\end{align}

\noindent\textbf{Step 4: The vector field $F_u$ is not conservative.} 
We show in this step that the vector field $F_u$ defined in \eqref{eq:def-F_u-main} is not a gradient field. 
In particular, we will show the existence of a rectangle $R \subset \mathbb{R}^2$ in the hidden space over which the path integral is negative. 
Set the comparator $u = (0, \frac{\pi}{6}) \in \mathcal{X}$ to show that $F_u$ is not a gradient field. 
Note that it is enough to consider a single comparator $u$ in order to show the linear regret bound. 
Denoting the coordinate functions of $F_u(z)$ by $P_u(z)$ and $Q_u(z)$ respectively, it is easy to verify at $z = (1,1)$ for instance 
that $\partial_{z_1} Q_u(z) - \partial_{z_2} P_u(z) < 0$.
By continuity of $(z_1, z_2) \mapsto \partial_{z_1} Q_u(z) - \partial_{z_2} P_u(z)$, 
there exists a rectangle $R := [a, a+ \alpha] \times [b, b+\beta] \subset q(\mathcal{X})$ for some $\alpha, \beta > 0$ and $a,b \in q(\mathcal{X})$ such that:
there exists a constant $\gamma > 0$ such that for all $z \in R$, $\partial_{z_1} Q_u(z) - \partial_{z_2} P_u(z) \leq - \gamma$\,. 
We denote by $A := (a,b), B:= (a+ \alpha, B), C:= (a+\alpha, b+ \beta)$ and $D:= (a, b + \beta)$ its corners and 
by $I_R$ the continuous oriented integral over the edges of the rectangle $R$ in the counter-clockwise orientation ($A \to B \to C \to D$). 
It follows then that: 
\begin{equation}
\label{eq:integral-I_R-main}
    I_R = \int_a^{a + \alpha} \int_b^{b+\beta} (\partial_{z_1} Q_u - \partial_{z_2} P_u)(x,y) dx dy \leq - \gamma \alpha \beta < 0\,.
\end{equation}

\noindent\textbf{Step 5: Definition of the adversarial sequence $(s_t)$.} 
Recall that we have set $s_t = - G(z_t) v_t$. 
We now choose the sequence $(v_t)$ in an adaptive adversarial way that makes the hidden iterates move counter-clockwise around the rectangle $R$. 
More precisely, we choose $v_t \in \{e_1, e_2, -e_1, -e_2 \}$ where $e_1 = (1,0)^{\top}, e_2 = (0,1)^{\top}$ in the following way, 
defined periodically for any $t$: $v_t = e_1$ along $AB$, $e_2$ along $BC$, $-e_1$ along $CD$ and $-e_2$ along $DA$. 
Let $n_1 := \lfloor \frac{\alpha}{\eta}\rfloor, n_2 :=  \lfloor \frac{\beta}{\eta}\rfloor$. 
Then the full cycle around the rectangle $R$ consists of $N = 2 n_1 + 2 n_2 = \frac{2(\alpha + \beta)}{\eta} + \mathcal{O}(1)$ steps. 
Then $h_t(z) = \langle s_t, z \rangle$ at round $t$ with $s_t = - G(z_t) v_t\,.$ 
This gives $\ell_t(x) =h_t(q(x)) = \langle s_t, q(x)\rangle$ and defines entirely the sequence of adversarial loss functions over which OGD is run. 

\noindent\textbf{Step 6: Riemann sum estimate of the path sum in regret in \eqref{eq:regret-as-path-sum} using the rectangle integral $I_R$.} 
Our goal in this step is to relate the path sum appearing in the regret decomposition in \eqref{eq:regret-as-path-sum} 
and the negative rectangle integral $I_R$ in \eqref{eq:rect-integral-neg-bound}. Specifically we show that: 
\begin{equation}
    \label{eq:path-sum-to-I_R1}
    \sum_{t=1}^N \langle F_u(z_t), z_{t+1} - z_t \rangle = I_R + \mathcal{O}(\eta)\,,
    \end{equation}
    where $N = 2(n_1 + n_2)$ is the number of steps in one cycle of the iterates around the rectangle $R$ 

\noindent\textbf{Step 7: Regret over one cycle of length $N$ around the rectangle $R$.} 
We show that one full cycle results in regret of order $\frac{1}{\eta}$. 
Since one cycle lasts $N = \Theta(\frac{1}{\eta})$ rounds, this gives a positive constant regret per round (on average).
Summing up \eqref{eq:reg-compon-decomp} over one cycle of $N$ steps yields: 
    \begin{equation}
    \label{eq:path-sum-to-I_R2}
    \sum_{t=1}^N \ell_t(x_t) - \ell_t(u) = - \frac{1}{\eta} \sum_{t=1}^N \langle F_u(z_t), z_{t+1} - z_t \rangle + \mathcal{O}(N \eta)
    = - \frac{1}{\eta} I_R + \mathcal{O}(1)\,, 
    \end{equation}
    where the first equality uses the fact that one cycle has $N = \Theta(\frac{1}{\eta})$ steps 
    and the second identity follows from using \eqref{eq:path-sum-to-I_R1}. 
    Since $I_R \leq - \gamma \alpha \beta < 0$, it follows that: 
    \begin{equation}
        \label{eq:1-cycle-lower-bound}
    \sum_{t=1}^N \ell_t(x_t) - \ell_t(u) \geq + \frac{\gamma \alpha \beta} {\eta} + \mathcal{O}(1) \geq \frac{c_0}{\eta}\,, 
    \end{equation}
    for some $c_0 > 0$ and for $\eta \in (0, \eta_0]$ for some $\eta_0 > 0$. 
    
    \noindent\textbf{Step 8: Linear regret lower bound by repeating cycles.} 
    We conclude the proof of Theorem~\ref{thm:linear-reg-lower-bound} by repeating the previous cycle over the time horizon $T$. 
    Let $K := \lfloor \frac{T}{N} \rfloor$ for $T \geq N$  be the number of complete cycles up to time $T$. 
    Since each cycle contributes at least $\frac{c_0}{\eta}$ regret, summing up these contributions over the $K$ cycles concludes the proof.   

\subsection{Proof of Theorem~\ref{thm:linear-reg-lower-bound}}
\label{appx:complete-proof-thm:linear-reg-lower-bound}

\textbf{Step 1: Definition of the parameterization $q$ and induced metric.} We start by defining the parameterization function $q$ to prove the result. 

Let $\mathcal{X} := [0,\log 3] \times [\frac \pi 6, \frac \pi 3] \subset \mathbb{R}^2.$ Define the parameterization function $q: \mathcal{X} \to \mathbb{R}^2$ for any $x = (x_1, x_2) \in \mathcal{X}$ as follows: 
\begin{equation*}
q(x_1,x_2) := e^{x_1} (\cos x_2, \sin x_2)\,.
\end{equation*}
Then it follows that the Jacobian of $q$ at any $x = (x_1, x_2) \in \mathcal{X} $ is given by: 
\begin{equation*}
J_q(x) = 
\begin{pmatrix}
\frac{\partial q_1(x)}{\partial x_1} & \frac{\partial q_1(x)}{\partial x_2} \\
\frac{\partial q_2(x)}{\partial x_1} & \frac{\partial q_2(x)}{\partial x_2}
\end{pmatrix} 
= e^{x_1} 
\begin{pmatrix} 
 \cos x_2 &  -  \sin x_2\\ 
 \sin x_2 &  \cos x_2
\end{pmatrix}
\,.
\end{equation*}
Therefore, we obtain: 
\begin{equation*}
J_q(x) J_q(x)^{\top} = e^{2x_1} I_2\,, 
\end{equation*}
where $I_2$ is the $2\times 2$ identity matrix. Since $x_1 \geq 0$, $e^{2x_1} \geq 1$ and hence $J_q(x) J_q(x)^{\top} \succeq I_2$. 
Moreover, setting $z := q(x)$, we have $\|z\| = \|q(x)\| = e^{x_1}$ and therefore: 
\begin{equation*}
J_q(x) J_q(x)^{\top} = \|z\|^2 I_2\,.
\end{equation*}
We now introduce additional notation which will be useful in the rest of the proof. Define for any $z \in q(\mathcal{X})$,  
\begin{align*}
M(z) &:= J_q(x) J_q(x)^{\top} = \|z\|^2 I_2\,,\\ 
G(z) &:= M(z)^{-1} = \frac{1}{\|z\|^2} I_2\,.
\end{align*}

\noindent\textbf{Step 2: OGD dynamics in the hidden space.} To construct our adversarial sequence of losses and give more insight about the idea of the proof, we first write the OGD dynamics in the hidden space $\Y := q(\mathcal{X}).$ 

We will consider linear functions $h_t: \Y \to \mathbb{R}$ defined for every $z \in \Y$ by $h_t(z) = \langle s_t, z \rangle$ where $(s_t)$ is an adversarial sequence that will be specified later in an adversarial way. It follows that for any $x \in \mathcal{X}$, 
\begin{equation*}
\ell_t(x) = h_t(q(x)) = \langle s_t, q(x) \rangle\,, \quad \nabla \ell_t(x) = J_q(x)^{\top} s_t. 
\end{equation*}

We first record the local hidden-space dynamics in the case where the Euclidean projection is inactive. For (projected) OGD on the sequence of function $(\ell_t)$, the update rule is:
\begin{equation*}
x_{t+1}
= \Pi_{\mathcal X}\bigl(x_t-\eta \nabla \ell_t(x_t) \bigr) = 
\Pi_{\mathcal X}\bigl(x_t-\eta J_q(x_t)^\top s_t\bigr).
\end{equation*}
Thus, on any round for which the pre-projection point $\widetilde x_{t+1}:=x_t-\eta J_q(x_t)^\top s_t$
belongs to $\mathcal X$, the projection is inactive and
\begin{equation}
\label{eq:ogd-update-x}
x_{t+1} = x_t - \eta \nabla \ell_t(x_t) = x_t - \eta J_q(x_t)^{\top} s_t\,.
\end{equation}
We will verify below (in step 5), after constructing the adversarial cycle, that the iterates remain
a positive distance away from the boundary $\partial\mathcal X$. Consequently, for sufficiently small
$\eta$, the pre-projection points remain in $\mathcal X$, so the projection is indeed inactive throughout the constructed trajectory.

Now defining the sequence $(z_t)$ by $z_t := q(x_t)$ for any $t$ in hidden space $\Y$, we obtain the following using Taylor's theorem with exact remainder: 
\begin{equation}
\label{eq:1-taylor}
z_{t+1} = q(x_{t+1}) = q(x_t) + J_q(x_t) (x_{t+1} - x_t) + r_t\,,
\end{equation}
where $r_t$ is a Taylor remainder which can be bounded as follows by global boundedness of the Hessian of $q$ (recall here that the space $\mathcal{X}$ is compact): 
\begin{equation*}
\|r_t\| = \mathcal{O }(\|x_{t+1} -x_t\|^2) = \mathcal{O }( \eta^2 \|\nabla \ell_t(x_t)\|^2) =  \mathcal{O }( \eta^2)\,, 
\end{equation*}
where the last step follows from uniform boundedness of the gradients of $f_t$ on the compact space $\mathcal{X}$ as $q$ is smooth and $(s_t)$ will be chosen as a bounded sequence. 

Plugging the OGD update rule of $(x_t)$ (see \eqref{eq:ogd-update-x}) in \eqref{eq:1-taylor}, we obtain: 
\begin{align*}
z_{t+1} &= z_t - \eta J_q(x_t) \nabla \ell_t(x_t) + r_t \nonumber\\
        &= z_t - \eta J_q(x_t) J_q(x_t)^{\top} s_t + r_t \nonumber\\
        &=  z_t - M(z_t) s_t + r_t\,. 
\end{align*}
To simplify the OGD dynamics in the hidden space, we choose $s_t = - G(z_t) v_t = - M(z_t)^{-1} v_t$ where $(v_t)$ is an adversarial sequence that will be chosen later. Given this choice, we obtain the following OGD dynamics in the hidden space: 
\begin{equation}
\label{eq:ogd-hidden-space}
z_{t+1} = z_t + \eta v_t + r_t\,,
\end{equation}
where $\|r_t\| = \mathcal{O}(\eta^2)$ as previously mentioned.\\

\noindent\textbf{Step 3: Regret as an approximate discrete path sum.} In this step, we write the regret under a specific path sum form for some field $F_u$ depending on a fixed comparator $u \in \mathcal{X}$. The main idea of the proof will then be to show that $F_u$ is not a gradient field and we will choose the sequence $(v_t)$ in an appropriate way to exploit this fact. Non-vanishing regret will be accumulated along the path due some cycling behavior of $F_u$. 

For a fixed comparator $u \in \mathcal{X}$, let $z_u := q(u).$ Then we have
\begin{align}
\ell_t(x_t) - \ell_t(u) &= h_t(q(x_t)) - h_t(q(u)) \nonumber\\ 
&= h_t(z_t) - h_t(z_u) \nonumber\\
&= \langle s_t, z_t - z_u \rangle  \nonumber\\
&=  - \langle G(z_t) v_t, z_t - z_u \rangle &&\hfill \text{($s_t = - G(z_t) v_t$)}   \nonumber\\
&= - \langle G(z_t) z_t - z_u, v_t  \rangle &&\hfill \text{(by symmetry of $G(z_t)$)}. \label{eq:3-fun-comp-gap}
\end{align}
Now define the vector field $F_u: \mathcal{Z} \to \mathbb{R}^2$ for any $z \in \mathcal{Z}$ by:
\begin{equation}
\label{eq:def-F_u}
F_u(z) := G(z) (z - z_u)\,.
\end{equation}
Using \eqref{eq:3-fun-comp-gap} together with the OGD dynamics \eqref{eq:ogd-hidden-space} in the hidden space, we obtain: 
\begin{align}
\ell_t(x_t) - \ell_t(u) &= - \langle F_u(z_t), v_t \rangle \nonumber\\
                  &= - \frac{1}{\eta} \langle F_u(z_t), z_{t+1} - z_t - r_t \rangle\nonumber\\ 
                  &= - \frac{1}{\eta} \langle F_u(z_t), z_{t+1} - z_t\rangle + \frac{1}{\eta}  \langle F_u(z_t), r_t \rangle\,.
\end{align}
Since $F_u$ is bounded and $\|r_t\| = \mathcal{O}(\eta^2)$, we have: 
\begin{equation*}
\left| \frac{1}{\eta}  \langle F_u(z_t), r_t \rangle \right| = \mathcal{O}\left(\frac{1}{\eta} \|r_t\|\right) = \mathcal{O}(\eta)\,.
\end{equation*}
Overall, we have 
\begin{equation}
\label{eq:reg-compon-decomp}
\ell_t(x_t) - \ell_t(u) = - \frac{1}{\eta} \langle F_u(z_t), z_{t+1} - z_t\rangle +  \mathcal{O}(\eta)\,.
\end{equation}
We have written regret as a path sum of the vector field $F_u$ up to an error of order $\mathcal{\eta T}$: 
\begin{align}
\label{eq:regret-as-path-sum}
R_T := \sum_{t=1}^T \ell_t(x_t) - \ell_t(u)  
= - \frac{1}{\eta} \sum_{t=1}^T  \langle F_u(z_t), z_{t+1} - z_t\rangle +  \mathcal{O}(\eta T)\,.
\end{align}

\noindent\textbf{Step 4: The vector field $F_u$ is not conservative.} We show in this step that the vector field $F_u$ defined in \eqref{eq:def-F_u} is not a gradient field. In particular, we will show the existence of a rectangle in the hidden space over which the path integral is negative. 

Set the comparator $u = (0, \frac{\pi}{6}) \in \mathcal{X}$ to show that $F_u$ is not a gradient field. Note that it will be enough for our purpose to consider a single comparator $u$ to show our linear regret bound. 

Then $z_u = q(u) = (\cos \frac{\pi}{6},\sin  \frac{\pi}{6}) = (\frac{\sqrt{3}}{2}, \frac 12)$. Recalling the definition of $F_u$ in \eqref{eq:def-F_u}, we have for any $z = (z_1, z_2) \in Q(\mathcal{X})$, 
\begin{equation*}
F_u(z) = G(z) (z - z_u) = \frac{1}{\|z\|^2} (z-z_u) = \frac{1}{z_1^2 + z_2^2} (z-z_u)\,.
\end{equation*}

Denote the coordinate functions of $F_u(z)$ as $P_u(z)$ and $Q_u(z)$. Then for any $z = (z_1, z_2) \in q(\mathcal{X})$, 
\begin{equation*}
F_u(z) = (P_u(z), Q_u(z)) = \left( \frac{z_1 - z_{u,1}}{z_1^2 + z_2^2}, \frac{z_2 - z_{u,2}}{z_1^2 + z_2^2}  \right)\,.
\end{equation*}
We now compute the difference of partial derivatives to obtain: 
\begin{equation*}
\partial_{z_1} Q_u(z) - \partial_{z_2} P_u(z) = \frac{2(z_1 z_{u,2} - z_2 z_{u_1})}{(z_1^2 + z_2^2)^2}\,.
\end{equation*}
At $z = (1, 1)$, we have: 
\begin{equation*}
\partial_{z_1} Q_u(z) - \partial_{z_2} P_u(z) = \frac 12 \left( \frac 12 - \frac{\sqrt{3}}{2}\right) < 0\,.
\end{equation*}

Since
$(1,1)=q(\log\sqrt 2,\pi/4)\in q(\operatorname{int}\mathcal X)$,
and since
$z\mapsto \partial_{z_1}Q_u(z)-\partial_{z_2}P_u(z)$
is continuous and strictly negative at $z=(1,1)$, there exist
\(\alpha,\beta>0\), \(a,b\in\mathbb R\), and \(\gamma>0\) such that: 
\begin{equation}
\label{eq:rec-R}
(1,1) \in R:=[a,a+\alpha]\times[b,b+\beta] \subset q(\operatorname{int}\mathcal X),
\end{equation}
where $R$ is a rectangle, and for all $z\in R$,
\begin{equation}
\label{eq:cross-partial-diff-neg}
\partial_{z_1}Q_u(z)-\partial_{z_2}P_u(z)\le -\gamma .
\end{equation}
Note here that we are making sure the rectangle $R$ is in $q(\operatorname{int}\mathcal X)$ is order to be able to ignore the projection in OGD. 
Moreover, by shrinking $R$ if necessary, we may assume that there exists
$\rho_z>0$ such that the compact set: 
\begin{equation}
\label{eq:def-Kz}
K_z:=\{z\in\mathbb R^2:\operatorname{dist}(z,R)\le \rho_z\}.
\end{equation}
satisfies $K_z \subset q(\operatorname{int}\mathcal X)$. 
This neighborhood is a buffer around the rectangle which ensures that the actual perturbed discrete trajectory $z_t$ (induced by the adversarial losses we will define) remains inside $q(\operatorname{int}\mathcal X)$, not just the ideal sequence which goes exactly around the rectangle without errors. This will be used  in steps 5 and 6 below.  

Let the corners of the rectangle $R$ be denoted as follows: 
\begin{align*}
D &:= (a, b + \beta)\,, && C:= (a+\alpha, b+ \beta)\,,\\
A &:= (a,b)\,,  && B:= (a+ \alpha, B)\,.
\end{align*}

We consider the counter-clockwise orientation: 
\begin{equation*}
A \to B \to C \to D\,.
\end{equation*}
Given this orientation, we define the following continuous edge integral: 
\begin{equation}
\label{eq:def-integral-I_R}
I_R := \int_A^B P_u(z) dz_1 + \int_B^C Q_u(z) dz_2 + \int_C^D P_u(z) dz_1 + \int_D^A Q_u(z) dz_2\,.
\end{equation}

The above integral can be rewritten as follows using differential calculus: 
\begin{equation}
\label{eq:integral-I_R}
    I_R = \int_a^{a + \alpha} \int_b^{b+\beta} (\partial_{z_1} Q_u - \partial_{z_2} P_u)(x,y) dx dy\,.
\end{equation}
The reader familiar with differential calculus may immediately see that the result holds. We provide an elementary proof for of this identity below for completeness:  
\begin{proof}
(of \eqref{eq:integral-I_R})
By definition of the continuous edge integral and the coordinate functions, we have: 
\begin{align}
I_R &= \int_a^{a + \alpha} P_u(x,b) dx + \int_b^{b + \beta} Q_u(a + \alpha, y) dy - \int_a^{a+ \alpha} P_u(x,b+\beta) dx - \int_b^{b+\beta} Q_u(a,y) dy \nonumber\\
&= \int_a^{a+\alpha} (P_u(x,b) - P_u(x,b+\beta))dx + \int_b^{b+\beta} (Q_u(a+\alpha,y) - Q_u(a,y))dy \nonumber\\
&= - \int_a^{a+\alpha} \int_b^{b+\beta} \partial_{z_2} P_u(x,y) dy dx + \int_b^{b+\beta} \int_a^{a+\alpha} \partial_{z_1}Q_u(x,y) dx dy\,,
\end{align} 
where the third equality follows from the one-dimensional fundamental theorem of calculus. 
\end{proof}

Using \eqref{eq:cross-partial-diff-neg}, we obtain that: 
\begin{equation}
\label{eq:rect-integral-neg-bound}
I_R \leq \int_a^{a + \alpha} \int_b^{b+\beta} (-\gamma) dx dy = - \gamma \alpha \beta < 0\,.
\end{equation}

\noindent\textbf{Step 5: Definition of the adversarial sequence $(s_t)$ via $(v_t)$.} We choose the sequence $(v_t)$ in an adaptive adversarial way that makes the hidden iterates move counter-clockwise around the rectangle $R$ defined in \eqref{eq:rec-R}. 

More precisely, we choose $v_t \in \{e_1, e_2, -e_1, -e_2 \}$ where $e_1 = (1,0)^{\top}, e_2 = (0,1)^{\top}$ in the following way, defined periodically for any $t$:
\begin{itemize}
\item along the edge $AB$: $v_t = e_1$\,,
\item along the edge $BC$: $v_t = e_2$\,, 
\item along the edge $CD$: $v_t = -e_1$\,, 
\item along the edge $DA$: $v_t = -e_2$\,. 
\end{itemize}

Let $n_1 := \lfloor \frac{\alpha}{\eta}\rfloor, n_2 :=  \lfloor \frac{\beta}{\eta}\rfloor$. Then the full cycle around the rectangle $R$ consists of $N = 2 n_1 + 2 n_2 = \frac{2(\alpha + \beta)}{\eta} + \mathcal{O}(1)$ steps. Then $h_t(z) = \langle s_t, z \rangle$ at round $t$ with $s_t = - G(z_t) v_t\,.$ This gives $\ell_t(x) =h_t(q(x)) = \langle s_t, q(x)\rangle$ and defines entirely the sequence of adversarial loss functions over which OGD is run. 

We next verify that, for this adversarial sequence and for sufficiently small $\eta$, the Euclidean projection in projected OGD is never active. This allows us to use the same local hidden-space expansion as in the unconstrained update, i.e., it remains to justify the claim made in Step 2 that the Euclidean projection is inactive along the constructed trajectory. Since
$K_z\subset q(\operatorname{int}\mathcal X)$,
the set $K_x:=q^{-1}(K_z)$ is a compact subset of $\operatorname{int}\mathcal X$, using continuity of $q^{-1}$. Therefore
$\rho_x:=\operatorname{dist}(K_x,\partial\mathcal X)>0.$
Moreover, on $K_z$, the adversarial vectors are uniformly bounded. Indeed,
\(v_t\in\{\pm e_1,\pm e_2\}\), \(s_t=-M(z_t)v_t\), and \(M\) is continuous on the compact set
$K_z$. Hence there exists $C_s<\infty$ such that
$\|s_t\|\le C_s$ whenever $z_t\in K_z$.
Since $J_q$ is continuous on the compact set $K_x$, there exists $C_J<\infty$ such that $\|J_q(x)^\top\|\le C_J$
for all $x\in K_x$. 
Thus, whenever $z_t\in K_z$ and $x_t=q^{-1}(z_t)$,
$
\|\eta J_q(x_t)^\top s_t\|
\le
\eta C_JC_s.
$
Choosing $\eta\le \rho_x/(2C_JC_s)$, we get
$\widetilde x_{t+1}=x_t-\eta J_q(x_t)^\top s_t\in \mathcal X.$
Therefore $\Pi_{\mathcal X}(\widetilde x_{t+1})=\widetilde x_{t+1}.$
So projected OGD coincides with the unconstrained update on every round for which $z_t\in K_z$.\\ 

\noindent\textbf{Step 6: Riemann sum estimate of the path sum in regret in \eqref{eq:regret-as-path-sum} using the rectangle integral $I_R$.} Our goal in this section is to relate the path sum appearing in the regret decomposition in \eqref{eq:regret-as-path-sum} and the negative rectangle integral $I_R$ in \eqref{eq:rect-integral-neg-bound}. For this we will rely on an edge integral approximation lemma. We present the result for the horizontal edge $AB$. Similar results hold for the other edges of the rectangle with minor modifications.  
Recall the OGD dynamics in the hidden space. For any $t \geq 0,$ 
\begin{equation*}
z_{t+1} = z_t + \eta v_t + r_t\,.
\end{equation*}

\begin{tcolorbox}[colframe=white!, top=2pt,left=2pt,right=2pt,bottom=2pt] 
\begin{lemma}[Horizontal edge approximation lemma]
\label{lem:horizontal-edge-approx}
Let $F = (P, Q)$ be $C^1$ on a compact neighborhood of the horizontal segment $E := \{ (x,b) : x \in  [a, a+\alpha] \}.$ Suppose that a discrete sequence $z_0, \dots, z_n$ for $n = \lfloor \frac{\alpha}{\eta} \rfloor$ satisfies for all $0 \leq k \leq n-1$, 
\begin{equation}
\label{eq:update-rule-zk-horizontal}
z_{k+1} = z_k + \eta e_1 + r_k\,,
\end{equation}
where $\|r_k\| = O(\eta^2)$ with $z_0 = (a,b)$. Then we have: 
\begin{equation*}
\sum_{k=0}^{n-1} \langle F(z_k), z_{k+1} - z_k \rangle = \int_a^{a+\alpha} P(x,b) dx + \mathcal{O}(\eta)\,.
\end{equation*}
\end{lemma}
\end{tcolorbox}

\begin{proof}
First, define the following auxiliary grid on the segment $[a, a+\alpha]$, 
\begin{equation*}
\bar{z}_k := (a + k \eta, b), \quad k= 0, \dots, n\,.
\end{equation*}
Then it follows that: 
\begin{equation*}
z_k = \bar{z}_k + \sum_{j=0}^{n-1} r_j\,.
\end{equation*}
and therefore we have for any $1 \leq k \leq n$, 
\begin{equation}
\label{eq:diff-zk-bar-zk}
\|z_k - \bar{z}_k \| = \left\| \sum_{j=0}^{k-1} r_j \right\| = \mathcal{O}(n \eta^2) = \mathcal{O}(\alpha \eta)\,,  
\end{equation}
where the last estimate follows from recalling that $n = \lfloor \frac{\alpha}{\eta}\rfloor\,.$ 
This shows that $z_k$ stays $\mathcal{O}(\eta)$-close to~$\bar{z}_k\,.$ 

Now we can estimate the path sum in \eqref{eq:regret-as-path-sum} as follows: 
\begin{align}
\label{eq:riemann-estimate-proof}
\sum_{k=0}^{n-1} \langle F(z_k), z_{k+1} - z_k \rangle 
&= \sum_{k=0}^{n-1} \langle F(z_k), \eta e_1 + r_k \rangle \nonumber\\
&= \eta \sum_{k=0}^{n-1} P(z_k) + \sum_{k=0}^{n-1} \langle F(z_k), r_k\rangle \nonumber\\
&= \eta \sum_{k=0}^{n-1} P(\bar{z}_k) + \eta \eta \sum_{k=0}^{n-1} P(z_k) - P(\bar{z}_k) + \sum_{k=0}^{n-1} \langle F(z_k), r_k\rangle\,,
\end{align}
where the first identity follows from using \eqref{eq:update-rule-zk-horizontal}, the second one follows from recalling that $P$ is the first coordinate function of $F$. 

We estimate each one of the three terms above separately. 

\noindent\textit{Last term in \eqref{eq:riemann-estimate-proof}.} As $F$ is bounded by continuity of $F$ on the cimpact neighborhood, we obtain: 
\begin{equation}
\label{eq:term3}
\sum_{k=0}^{n-1} \langle F(z_k), r_k \rangle = \mathcal{O}\left(n \cdot \max_{0 \leq k \leq n-1} \|r_k\|\right) = \mathcal{O}(n \eta^2) = \mathcal{O}(\eta)\,,
\end{equation}
where the last step uses again the fact that $n = \lfloor \frac{\alpha}{\eta}\rfloor.$ 

\noindent\textit{Second term in \eqref{eq:riemann-estimate-proof}.} Using Lipschitzness of $P$ and recalling that $\max_{0 \leq k \leq n-1} \|z_k - \bar{z}_k\| = \mathcal{O}(\eta)$ proved in \eqref{eq:diff-zk-bar-zk}, we have $P(z_k) - P(\bar{z}_k) = \mathcal{O}(\eta)$ and it follows that: 
\begin{equation}
\label{eq:term2}
\eta \sum_{k=0}^{n-1} P(z_k) - P(\bar{z}_k) = \mathcal{O}(n \eta^2) = \mathcal{O}(\eta)\,.
\end{equation}

\noindent\textit{First term in \eqref{eq:riemann-estimate-proof}.} Denoting $\bar{x}_k = a + k \eta$ for $0 \leq k \leq n-1$, we can write: 
\begin{align}
\label{eq:term1}
\eta \sum_{k=0}^{n-1} P(\bar{z}_k) 
&= \sum_{k=0}^{n-1} \int_{\bar{x}_k}^{\bar{x}_{k+1}} P(\bar{z}_k) dx \nonumber\\
&= \sum_{k=0}^{n-1} \int_{\bar{x}_k}^{\bar{x}_{k+1}} P(x,b) dx 
+ \sum_{k=0}^{n-1} \int_{\bar{x}_k}^{\bar{x}_{k+1}} P(\bar{z}_k) -  P(x,b) dx \nonumber\\
&= \int_0^{a+n \eta} P(x,b) dx + R_n\,, 
\end{align}
where the residual $R_n$ is defined as follows: 
\begin{equation*}
R_n := \sum_{k=0}^{n-1} \int_{\bar{x}_k}^{\bar{x}_{k+1}} P(\bar{z}_k) - P(x,b) dx\,.
\end{equation*}

Using again Lipschitzness of $P$, the above residual can be bounded as follows: 
\begin{align}
\label{eq:term1-Rn-bound}
|R_n| &\leq \sum_{k=0}^{n-1} \int_{\bar{x}_k}^{\bar{x}_{k+1}} |P(\bar{z}_k) - P(x,b)| \,dx \nonumber\\ 
&= \mathcal{O}\left( \sum_{k=0}^{n-1} \int_{\bar{x}_k}^{\bar{x}_{k+1}} |\bar{x}_k - x| \,dx \right) \nonumber\\
&=  \mathcal{O}\left( n \cdot  \max_{0 \leq k \leq n-1} |\bar{x}_{k+1} - \bar{x}_k|^2 \right) \nonumber\\
&= \mathcal{O}(n \eta^2) \nonumber\\
&= \mathcal{O}(\eta)\,. 
\end{align}
Now, observe that: 
\begin{equation*}
\int_a^{a+\alpha} P(x,b) dx = \int_a^{a+n \eta} P(x,b) dx + \int_{a+n \eta}^{a+\alpha} P(x,b) dx\,.
\end{equation*}
Hence, since $0 \leq \alpha - n \eta \leq \eta$ by definition of $n = \lfloor \frac{\alpha}{\eta} \rfloor$ and since $P(x,b)$ is bounded over the compact neighborhood of the horizontal segment, we obtain: 
\begin{equation}
\label{eq:segment-discrete-residual}
\int_a^{a+\alpha} P(x,b) dx = \int_a^{a+n \eta} P(x,b) dx + \mathcal{O}(\eta)\,.
\end{equation}
Combining all the above estimates (\eqref{eq:term3}, \eqref{eq:term2}, \eqref{eq:term1}, \eqref{eq:term1-Rn-bound}  and \eqref{eq:segment-discrete-residual}) in the main term \eqref{eq:riemann-estimate-proof}, we obtain the following estimate: 
\begin{equation*}
\sum_{k=0}^{n-1} \langle F(z_k), z_{k+1} - z_k \rangle = \int_a^{a+\alpha} P(x,b) dx + \mathcal{O}(\eta)\,,
\end{equation*}
concluding the proof of Lemma~\ref{lem:horizontal-edge-approx}. 
\end{proof}

The previous estimates (Lemma~\ref{lem:horizontal-edge-approx}) also imply that the actual discrete trajectory $(z_t)$ remains in~$K_z$ as defined in~\eqref{eq:def-Kz}.
Indeed, along each edge the one-step error is $\mathcal{O}(\eta^2)$, and each edge lasts $\mathcal{O}(1/\eta)$ rounds, so the accumulated deviation from the intended rectangular path is
$\mathcal{O}(\eta)$. Choosing $\eta$ sufficiently small so that this deviation is at most
$\rho_z$, the iterates remain in $K_z$ throughout the cycle. Consequently, by the conditional projection-inactivity argument above (see end of step 5), the projection is inactive at every round of the cycle.\\

\noindent\textbf{Step 7: Regret over one cycle of length $N$ around the rectangle $R$.} Applying Lemma~\ref{lem:horizontal-edge-approx} to the edge AB, and similarly to BC, CD, DA (which follow using similar proofs with minor adjustments, e.g. replacing $e_1$ by $-e_1$, $e_2$ and $-e_2$ and the coordinate function $P_u$ by $Q_u$), we sum all the estimates to obtain: 
\begin{equation*}
\sum_{t=1}^N \langle F_u(z_t), z_{t+1} - z_t \rangle = I_R + \mathcal{O}(\eta)\,,
\end{equation*}
where $N = 2(n_1 + n_2)$ is the number of steps in one cycle of the iterates around the rectangle $R = ABCD.$ 

Summing up \eqref{eq:reg-compon-decomp} over one cycle of $N$ steps yields: 
\begin{equation}
\label{eq:path-sum-to-I_R}
\sum_{t=1}^N \ell_t(x_t) - \ell_t(u) = - \frac{1}{\eta} \sum_{t=1}^N \langle F_u(z_t), z_{t+1} - z_t \rangle + \mathcal{O}(N \eta)\,.
\end{equation}

Since one cycle has $N = \Theta(\frac{1}{\eta})$ steps, we have $N \eta = \mathcal{O}(1)$. Combining this with \eqref{eq:path-sum-to-I_R}, we obtain: 
\begin{equation*}
\sum_{t=1}^N \ell_t(x_t) - \ell_t(u) = - \frac{1}{\eta} I_R + \mathcal{O}(1)\,.
\end{equation*}
Recalling from \eqref{eq:rect-integral-neg-bound} that we have proved that $I_R \leq - \gamma \alpha \beta < 0$, it follows that: 
\begin{equation*}
\sum_{t=1}^N \ell_t(x_t) - \ell_t(u) \geq + \frac{\gamma \alpha \beta} {\eta} + \mathcal{O}(1)\,.
\end{equation*}

As $\frac{\gamma \alpha \beta}{\eta} \to +\infty$ when $\eta \to 0$, there exists $\eta_0 > 0$ and $c_0 > 0$ such that for all $\eta \in (0, \eta_0]$,  
\begin{equation}
\label{eq:per-cycle-regret-bound}
\sum_{t=1}^N \ell_t(x_t) - \ell_t(u) \geq \frac{c_0}{\eta}\,. 
\end{equation} 
Therefore, we have shown that one full cycle results in regret of order $\frac{1}{\eta}$. Since one cycle lasts $N = \Theta(\frac{1}{\eta})$ rounds, this gives a positive constant regret per round (on average).\\

\noindent\textbf{Step 8: Linear regret lower bound by repeating cycles.} We conclude the proof of Theorem~\ref{thm:linear-reg-lower-bound} by repeating the previous cycle over the time horizon $T$. 

Let $K := \lfloor \frac{T}{N} \rfloor$ for $T \geq N$  be the number of complete cycles up to time $T$. We know from \eqref{eq:per-cycle-regret-bound} in the previous step that each cycle contributes at least $\frac{c_0}{\eta}$ regret. Thus, we have: 
\begin{equation*}
\sum_{t=1}^T \ell_t(x_t) - \ell_t(u) \geq K \frac{c_0}{\eta} - \mathcal{O}(1)\,.
\end{equation*}
Since $N \leq \frac{c_2}{\eta}$ for some constant $c_2 > 0$, we have: 
\begin{equation*}
K \geq \frac{T}{N} -1 \geq \eta \frac{T}{c_2} - 1\,.
\end{equation*}
Substituting this in the previous lower bound yields: 
\begin{equation*}
    \sum_{t=1}^T \ell_t(x_t) - \ell_t(u) \geq \left( \frac{\eta T}{c_2} - 1 \right) \frac{c_0}{\eta} - \mathcal{O}(1) = \frac{c_0}{c_2} T - \frac{c_0}{\eta} - \mathcal{O}(1) = \frac{c_0}{c_2} T - \mathcal{O}(1)\,, 
\end{equation*}
which concludes the proof of Theorem~\ref{thm:linear-reg-lower-bound}. 

\section{Proof of Theorem~\ref{thm:bandit}: Expected Regret Bound in the Bandit Setting}
\label{appx:proofs-bandit}

Recall first that the regret of bandit OGD (Algorithm~\ref{alg4}) is defined for any time horizon $T \geq 1$ by: 
\begin{equation}
\label{eq:regret-bogd}
R_T := \sum_{t=1}^T \ell_t(\hat{x}_t) - \min_{x \in \X} \sum_{t=1}^T \ell_t(x)\, 
= \sum_{t=1}^T h_t(\hat{z}_t) - \min_{z \in \Y} \sum_{t=1}^T h_t(z)\,,
\end{equation}
where the last equality follows from using the notation $\hat{z}_t := q(\hat{x}_t)$ for any $t \geq 1$ and recalling that $\ell_t = h_t \circ q$ with $q$ being a diffeomorphism. 

To prove Theorem~\ref{thm:bandit}, we first decompose the expected regret into four different error terms.

\begin{tcolorbox}[colframe=white!, top=2pt,left=2pt,right=2pt,bottom=2pt] 
\begin{lemma}[Expected regret bound decomposition]
\label{lem:expected-reg-decomp}
Fix a time horizon $T \geq 1$.  Let $z^{\star} \in \argmin_{z \in \Y} \sum_{t=1}^T h_t(z), z^{\star}_{\delta} \in \argmin_{z \in \Y_{\delta}} \sum_{t=1}^T h_t(z)\,$ and let $\mathcal{F}_{t} := \sigma(\zeta_1, \dots, \zeta_t)$ be the filtration of $\sigma$-algebras induced by the random variables $(\zeta_t)$ used in bandit OGD. Then the expected regret $\mathbb{E}[R_T]$ of bandit OGD (Algorithm~\ref{alg4}) can be upper-bounded as follows: 
\begin{equation*}
\mathbb{E}[R_T] \leq R_T^{(1)} +  R_T^{(2)} + R_T^{(3)} + R_T^{(4)}\,,
\end{equation*}
where each error term in the right-hand side above is defined as follows: 
\begin{align}
R_T^{(1)} &:= \mathbb{E}\left[\sum_{t=1}^T \langle \tilde{g}_t, z_t - z^{\star}_{\delta} \rangle \right]\,,  &&\hfill \text{(perturbed OMD regret for stochastic linearized losses)} \label{eq:RT1}\\
R_T^{(2)} &:= \sum_{t=1}^T h_t(z^{\star}_{\delta}) - h_t(z^{\star})\,,  &&\hfill \text{(domain shrinkage error)}\label{eq:RT2}\\
R_T^{(3)} &:= \mathbb{E}\left[\sum_{t=1}^T \langle b_t, z^{\star}_{\delta} - z_t \rangle \right]\,,  &&\hfill \text{(cumulative gradient estimation bias)}\label{eq:RT3}\\
b_t &:= \mathbb{E}[\tilde{g}_t |\mathcal{F}_{t-1}] - \nabla h_t(z_t)\,,  &&\hfill \text{(gradient estimation bias)}\\
R_T^{(4)} &:= \mathbb{E}\left[ \sum_{t=1}^T h_t(\hat{z}_t) - h_t(z_t)\right]  &&\hfill \text{(smoothing error)} \label{eq:RT4}\,.
\end{align}
\end{lemma}
\end{tcolorbox}

\begin{proof}
Using the regret definition~\eqref{eq:regret-bogd}, we have the following decomposition: 
\begin{equation}
\label{eq:regret-decomp1}
\mathbb{E}[R_T] = \mathbb{E}\left[\sum_{t=1}^T h_t(\hat{z}_t) - h_t(z^{\star}_{\delta})\right] 
+ \sum_{t=1}^T h_t(z^{\star}_{\delta}) - h_t(z^{\star})
= \mathbb{E}\left[\sum_{t=1}^T h_t(\hat{z}_t) - h_t(z^{\star}_{\delta})\right] + R_T^{(2)}\,.
\end{equation}
To control regret against the shrunk comparator $z_{\delta}^{\star}$, we start by decomposing the instantaneous regret. Since $z_t$ is $\mathcal{F}_{t-1}$-measurable, we have: 
\begin{equation}
\label{eq:1-reg-decomp}
\mathbb{E}[ h_t(\hat{z}_t) - h_t(z^{\star}_{\delta}) |\mathcal{F}_{t-1}] 
= \mathbb{E}[ h_t(\hat{z}_t) - h_t(z_t) |\mathcal{F}_{t-1}] 
+ h_t(z_t) -  h_t(z_{\delta}^{\star})\,.
\end{equation}
Then by convexity of $h_t$, we have: 
\begin{equation}
\label{eq:2-reg-decomp}
h_t(z_t) -  h_t(z_{\delta}^{\star}) \leq \langle \nabla h_t(z_t), z_t- z_{\delta}^{\star} \rangle\,.
\end{equation}
Recall now that we denote by $b_t$ the bias due to one-point smoothing defined by: 
\begin{equation*}
b_t = \mathbb{E}[\tilde{g}_t|\mathcal{F}_{t-1}] - \nabla h_t(z_t)\,.
\end{equation*}
Note here that $\tilde{g}_t := J_q(x_t)^{-\top} g_t$ is is not computed and just used here in the analysis whereas $g_t$ is the one-point estimator actually computed in BOGD. 
It follows then that: 
\begin{equation}
\label{eq:3-reg-decomp}
\langle \nabla h_t(z_t), z_t - z_{\delta}^{\star} \rangle 
= \langle \mathbb{E}[\tilde{g}_t|\mathcal{F}_{t-1}], z_t - z_{\delta}^{\star}\rangle - \langle b_t, z_t - z_{\delta}^{\star}\rangle\,.
\end{equation}
Combining \eqref{eq:1-reg-decomp}, \eqref{eq:2-reg-decomp} and \eqref{eq:3-reg-decomp}, taking expectations and noting that $z_t - z_{\delta}^{\star}$ is $\mathcal{F}_{t-1}$-measurable, we obtain: 
\begin{equation*}
\mathbb{E}[h_t(z_t) -  h_t(z_{\delta}^{\star})] 
\leq  \mathbb{E}[\langle \tilde{g}_t, z_t - z_{\delta}^{\star}\rangle] - \mathbb{E}[\langle b_t, z_t - z_{\delta}^{\star}\rangle] + \mathbb{E}[h_t(\hat{z}_t) - h_t(z_t)]\,.
\end{equation*}
Summing the above inequality over $t$ yields: 
\begin{align}
\label{eq:4-reg-decomp}
\mathbb{E}\left[ \sum_{t=1}^T h_t(z_t) -  h_t(z_{\delta}^{\star}) \right] 
&\leq  \mathbb{E}\left[\sum_{t=1}^T \langle \tilde{g}_t, z_t - z_{\delta}^{\star}\rangle \right] + \mathbb{E}\left[\sum_{t=1}^T \langle b_t, z^{\star}_{\delta} - z_t \rangle \right] + \mathbb{E}\left[ \sum_{t=1}^T h_t(\hat{z}_t) - h_t(z_t)\right]\nonumber\\ 
&= R_T^{(1)} + R_T^{(3)} + R_T^{(4)}. 
\end{align}
Finally, combining \eqref{eq:4-reg-decomp} with \eqref{eq:regret-decomp1} gives the desired bound: 
\begin{equation*}
\mathbb{E}[R_T] \leq R_T^{(1)} + R_T^{(2)} + R_T^{(3)} + R_T^{(4)}\,.
\end{equation*}
\end{proof}

Given the regret decomposition of Lemma~\ref{lem:expected-reg-decomp}, we control each of the error terms $R_T^{(i)}, i \in \{ 1,2,3,4\}$ defined in \eqref{eq:RT1} to \eqref{eq:RT4} respectively, in a dedicated lemma.   

\begin{tcolorbox}[colframe=white!, top=2pt,left=2pt,right=2pt,bottom=2pt] 
\begin{lemma}[Domain shrinkage error]
\label{lem:domain-shrinkage-error}
Let $T \geq 1$, let $z^{\star} \in \argmin_{z \in \Y} \sum_{t=1}^T h_t(z)$, and $z^{\star}_{\delta} \in \argmin_{z \in \Y_{\delta}} \sum_{t=1}^T h_t(z)\,.$ Then we have: 
\begin{equation*}
R_T^{(2)} = \sum_{t=1}^T h_t(z_{\delta}^{\star}) - h_t(z^{\star}) \leq \delta G_FDT\,,
\end{equation*}
where we recall that $\delta$ is the smoothing parameter in the one-point zeroth-order gradient estimator, $G_F$ is the uniform Lipschitz constant of $\ell_t$ and $D$ is the finite diameter of the compact set $\X$. 
\end{lemma}
\end{tcolorbox}

\begin{proof}
Since $\Y_{\delta} = q(\X_{\delta})$, $q$ is a diffeomorphism,  $\X_{\delta} = (1-\delta) \X$ and $\ell_t = h_t \circ q$, the comparator term over the shrunk set $\X_{\delta}$ can be rewritten as follows: 
\begin{equation}
\label{eq:min-eq}
\min_{z \in \Y_{\delta}} \sum_{t=1}^T h_t(z) = \min_{x \in \X_{\delta}} \sum_{t=1}^T \ell_t(x) = \min_{x \in \X} \sum_{t=1}^T \ell_t((1-\delta) x)\,.
\end{equation}
Then for any $x \in \X,$ we have: 
\begin{equation*}
\sum_{t=1}^T \ell_t((1-\delta)x)
\leq \sum_{t=1}^T |\ell_t((1-\delta)x) - \ell_t(x)| + \sum_{t=1}^T \ell_t(x) \leq \delta G_F DT + \sum_{t=1}^T h_t(q(x))\,,
\end{equation*}
where the last inequality uses uniform $G_F$-Lipschitzness of the loss functions $\ell_t$ and boundedness of the set $\X$ which has finite diameter $D$. 

Taking the minimum over $x \in \X$, using \eqref{eq:min-eq} and recalling that $z^{\star} \in \argmin_{z \in \Y} \sum_{t=1}^T h_t(z)$ and $z^{\star}_{\delta} \in \argmin_{z \in \Y_{\delta}} \sum_{t=1}^T h_t(z)$ yields the desired inequality: 
\begin{equation*}
R_T^{(2)} = \sum_{t=1}^T h_t(z_{\delta}^{\star}) - h_t(z^{\star}) \leq \delta G_FDT\,.
\end{equation*}
\end{proof}

\begin{tcolorbox}[colframe=white!,
top=2pt,left=2pt,right=2pt,bottom=2pt]
\begin{lemma}[Gradient estimation bias]
\label{lem:grad-estimation-bias}
For any $T \geq 1$, we have: 
\begin{equation*} 
|R_T^{(3)}| = \left| \mathbb{E}\left[ \sum_{t=1}^T \langle b_t, z_{\delta}^{\star} - z_t  \rangle \right]\right|   \leq \delta G^2 H DT\,,
\end{equation*}
where we recall that $H$ is the uniform smoothness constant of the loss $\ell_t.$
\end{lemma}
\end{tcolorbox}

\begin{proof}
First, using the Cauchy-Schwartz inequality, it follows that: 
\begin{equation}
\label{eq:ics-bias}
|R_T^{(3)}| = \left| \mathbb{E}\left[ \sum_{t=1}^T \langle b_t, z_{\delta}^{\star} - z_t  \rangle \right]\right|
\leq \mathbb{E}\left[ \sum_{t=1}^T \|b_t\| \cdot \|z_t - z_{\delta}^{\star}\| \right]\,. 
\end{equation}

We control each one of the terms in the right-hand side separately.

Since $z_t = q(x_t)$ and $z_{\delta}^{\star} = q(x_{\delta}^{\star})$ for some  $x_t \in \X$ and $x_{\delta}^{\star} \in \X_{\delta} \subset \X$, we have: 
\begin{equation}
\label{eq:bias-proof2}
\|z_t - z_{\delta}^{\star}\| = \|q(x_t) - q(x_{\delta}^{\star})\| \leq G \|x_t - x_{\delta}^{\star}\| \leq GD\,,
\end{equation}
where the first inequality follows from using $G$-Lipschitzness of $q$ (Assumption~\ref{as:reg-R-reparam-q-smoothness}) and the second one uses the definition of the diameter of $\X$ supposed to be finite.  

We now control the norm of the bias $\|b_t\|$ for any $t$. Define for any $x \in \X$ the smoothed loss: 
\begin{equation*}
\bar{\ell}_t(x) := \mathbb{E}_{v \sim \mathcal{U}(\mathbb{B})}[\ell_t(x + \delta v)]\,
\end{equation*}
where $v$ is a random variable with uniform distribution $\mathcal{U}(\mathbb{B})$ on the unit ball $\mathbb{B}$. It is known from \citet[Lemma~1]{flaxman-kalai-mcmahan05} that the conditional expectation of the spherical estimator $g_t$ is equal to the gradient of the ball-smoothed loss: 
\begin{equation*}
\mathbb{E}[g_t |\mathcal{F}_{t-1}] = \nabla \bar{\ell}_t(x_t)\,.
\end{equation*}
The (hidden-space ghost) gradient estimator of $\nabla h_t(z_t)$ is 
\begin{equation*}
\tilde{g}_t = J_q(x_t)^{-\top} g_t\,.
\end{equation*}
(Recall that $\nabla \ell_t(x_t) = J_q(x_t)^{\top} \nabla h_t(z_t)$ by the chain rule.)
Since $x_t$ is $\mathcal{F}_{t-1}$-measurable, it follows that: 
\begin{equation}
\label{eq:expectation-tilde-g}
\mathbb{E}[\tilde{g}_t|\mathcal{F}_{t-1}] = J_q(x_t)^{-\top} \mathbb{E}[g_t|\mathcal{F}_{t-1}] = J_q(x_t)^{-\top} \nabla \bar{\ell}_t(x_t)\,.
\end{equation}
Moreover, we have by the chain rule $\nabla \ell_t(x_t) = J_q(x_t)^{\top} \nabla h_t(z_t)$ which implies that $\nabla h_t(z_t)= J_q(x_t)^{-\top} \nabla \ell_t(x_t).$ Therefore, using this identity together with \eqref{eq:expectation-tilde-g}, we can express the bias $b_t$ as follows: 
\begin{align*}
b_t &= \mathbb{E}[\tilde{g}_t|\mathcal{F}_{t-1}] - \nabla h_t(z_t)\nonumber\\  
&= J_q(x_t)^{-\top} (\nabla \bar{\ell}_t(x_t) - \nabla \ell_t(x_t)) \nonumber\\ 
&= J_q(x_t)^{-\top} \mathbb{E}_{v \sim \mathcal{U}(\mathbb{B})}[\nabla \ell_t(x_t + \delta v) - \nabla \ell_t(x_t) ]\,.
\end{align*}
Taking the norm, using boundedness of the Jacobian of $q$ (Assumption~\ref{as:reg-R-reparam-q-smoothness}) and smoothness of the loss function $\ell_t$ (Assumption~\ref{as:smoothness-bandit}), we have: 
\begin{equation}
\label{eq:bias-proof3}
\|b_t\|_2 \leq \|J_q(x_t)^{-\top}\| \mathbb{E}_{v \sim \mathcal{U}(\mathbb{B})}[\|\nabla \ell_t(x_t + \delta v) - \nabla \ell_t(x_t) \|] \leq \delta G H \|v\|_2 = \delta GH\,. 
\end{equation}
Combining \eqref{eq:bias-proof3} and \eqref{eq:bias-proof2} in \eqref{eq:ics-bias}, we obtain the desired inequality: 
\begin{equation*}
|R_T^{(3)}| \leq \delta G^2 H D T\,.
\end{equation*}
\end{proof}

\begin{tcolorbox}[colframe=white!,
top=2pt,left=2pt,right=2pt,bottom=2pt]
\begin{lemma}[Smoothing error]
\label{lem:smoothing-error}
For any $T \geq 1$, we have: 
\begin{equation*}
|R_T^{(4)}| \leq \left| \mathbb{E}\left[ \sum_{t=1}^T h_t(\hat{z}_t) - h_t(z_t) \right] \right| \leq \delta G_F T\,.
\end{equation*}
\end{lemma}
\end{tcolorbox}

\begin{proof}
The desired bound follows from the following: 
\begin{align}
|R_T^{(4)}| &= \left| \mathbb{E}\left[ \sum_{t=1}^T h_t(\hat{z}_t) - h_t(z_t) \right] \right|\nonumber\\
&= \left| \mathbb{E}\left[ \sum_{t=1}^T \ell_t(\hat{x}_t) - \ell_t(x_t) \right] \right|&&\hfill \text{($\ell_t = h_t \circ q$, $\hat{z}_t = q(\hat{x}_t), z_t = q(x_t)$)}\nonumber\\
&\leq G_F \mathbb{E}\left[ \sum_{t=1}^T \|\hat{x}_t - x_t\| \right]
&&\hfill \text{(by $G_F$-Lipschitzness of $\ell_t$ (Assumption~\ref{as:boundedness-grad-bregman}))}\nonumber\\
&= G_F \mathbb{E}\left[ \sum_{t=1}^T \delta \|\zeta_t\| \right] &&\hfill \text{($\hat{x}_t = x_t + \delta \zeta_t$ in BOGD (Algorithm~\ref{alg4}))}\nonumber\\
&= \delta G_F T&&\hfill \text{($\|\zeta_t\| = 1$)}\,.
\end{align}
\end{proof}

\begin{tcolorbox}[colframe=white!,
top=2pt,left=2pt,right=2pt,bottom=2pt]
\begin{lemma}[Perturbed ghost OMD regret for random linear losses]
\label{lem:perturbed-omd-reg-random-lin-loss}
Let $z_1 \in \Y_{\delta}.$ Suppose the sequence $(z_t)$ is defined as: 
\begin{equation*}
z_{t+1} = r_{t+1} + y_{t+1}\,, \quad y_{t+1} = \argmin_{y \in \Y_{\delta}} \{ \tilde{g}_t^{\top} (y - z_t) + \frac{1}{\eta} D_R(y||z_t)\}\,.
\end{equation*}
Then for every $z_{\delta}^{\star} \in \Y_{\delta},$
\begin{equation*}
\sum_{t=1}^T \langle \tilde{g}_t, z_t - z_{\delta}^{\star} \rangle 
\leq \frac{D_R(z_{\delta}^{\star}||z_1)}{\eta} + \frac{\eta}{2} \sum_{t=1}^T \|\tilde{g}_t\|_2^2 + \frac{G}{\eta} \sum_{t=1}^T \|r_{t+1}\|_2\,.
\end{equation*}
If in addition $\|\tilde{g}_t\| \leq \tilde{G}_F$ and $\|r_{t+1}\| \leq \hat{C}_{\eta}$, then: 
\begin{equation}
\label{eq:R1T-bandit}
R_T^{(1)} = \mathbb{E}\left[ \sum_{t=1}^T \langle \tilde{g}_t, z_t - z_{\delta}^{\star} \rangle\right] \leq \frac{D_1}{\eta} + \frac{\eta \tilde{G}_F^2T}{2} + \frac{\hat{C}_{\eta} G T}{\eta}\,.
\end{equation}
\end{lemma}
\end{tcolorbox}

\begin{proof}
First, we prove the following inequality which follows from using the first-order optimality condition (from the ghost OMD update) together with Bregman three-point identity. For any comparator $z_{\delta}^{\star} \in \Y_{\delta}$, we have: 
\begin{equation}
\label{eq:omd-ineq}
\langle \tilde{g}_t, z_t - z_{\delta}^{\star}\rangle  
\leq \frac{D_R(z_{\delta}^{\star}||z_t) - D_R(z_{\delta}^{\star}||y_{t+1})}{\eta} + \frac{\eta}{2}\|g_t\|_2^2\,.
\end{equation}

We now provide the proof of \eqref{eq:omd-ineq}. By definition of $y_{t+1}$, using the first-order optimality condition gives for every $z_{\delta}^{\star} \in \Y_{\delta}$, 
\begin{equation*}
\left\langle \tilde{g}_t + \frac{1}{\eta} (\nabla R(y_{t+1}) - \nabla R(z_t)), z_{\delta}^{\star} - y_{t+1} \right\rangle \geq 0\,.
\end{equation*}
Rearranging this inequality, we obtain: 
\begin{equation}
\label{eq:proof-omd1}
\langle \tilde{g}_t, y_{t+1} - z_{\delta}^{\star} \rangle 
\leq \frac{1}{\eta} \langle \nabla R(y_{t+1}) - \nabla R(z_t), z_{\delta}^{\star} - y_{t+1} \rangle\,.
\end{equation}
Using the three-point identity, we have: 
\begin{equation}
\label{eq:proof-omd2}
\langle \nabla R(y_{t+1}) - \nabla R(z_t), z_{\delta}^{\star} - y_{t+1} \rangle = D_R(z_{\delta}^{\star}||z_t) - D_R(z_{\delta}^{\star}||y_{t+1}) - D_R(y_{t+1}||z_t)\,.
\end{equation}
Combining \eqref{eq:proof-omd1} and \eqref{eq:proof-omd2} yields: 
\begin{equation*}
\langle \tilde{g}_t, y_{t+1} - z_{\delta}^{\star} \rangle 
\leq \frac{1}{\eta}\left[ D_R(z_{\delta}^{\star}||z_t) - D_R(z_{\delta}^{\star}||y_{t+1}) - D_R(y_{t+1}||z_t) \right]\,.
\end{equation*}
Using the above inequality, we obtain: 
\begin{align*}
\langle \tilde{g}_t, z_t - z_{\delta}^{\star} \rangle 
&= \langle \tilde{g}_t, z_t - y_{t+1} \rangle 
+ \langle \tilde{g}_t, y_{t+1} - z_{\delta}^{\star} \rangle \nonumber\\ 
&\leq \langle \tilde{g}_t, z_t - y_{t+1} \rangle 
+ \frac{1}{\eta}\left[ D_R(z_{\delta}^{\star}||z_t) - D_R(z_{\delta}^{\star}||y_{t+1}) - D_R(y_{t+1}||z_t) \right] \nonumber\\
&= \langle \tilde{g}_t, z_t - y_{t+1} \rangle - \frac{1}{\eta}D_R(y_{t+1}||z_t) 
+ \frac{1}{\eta}\left[ D_R(z_{\delta}^{\star}||z_t) - D_R(z_{\delta}^{\star}||y_{t+1}) \right]\,.
\end{align*}
Using $1$-strong convexity of $R$ and Cauchy-Schwarz inequality yields: 
\begin{equation*}
\langle \tilde{g}_t, z_t - y_{t+1} \rangle - \frac{1}{\eta}D_R(y_{t+1}||z_t)  \leq \|\tilde{g}_t\|_2 \cdot \|z_t - y_{t+1}\|_2 - \frac{1}{\eta} \|z_t - y_{t+1}\|_2^2\,. 
\end{equation*}
Note now that for any $a \geq 0, r \geq 0, ar - \frac{1}{2\eta} r^2 \leq \frac{\eta}{2}a^2\,.$ Hence with $a = \|\tilde{g}_t\|_2$ and $r = \|z_t - y_{t+1}\|$, we obtain: 
\begin{equation*}
\langle \tilde{g}_t, z_t - y_{t+1} \rangle - \frac{1}{\eta}D_R(y_{t+1}||z_t)  \leq \frac{\eta}{2} \|\tilde{g}_t\|_2^2\,.
\end{equation*}
We have proved that: 
\begin{equation*}
\langle \tilde{g}_t, z_t - z_{\delta}^{\star} \rangle 
\leq \frac{D_R(z_{\delta}^{\star}||z_t) - D_R(z_{\delta}^{\star}||y_{t+1})}{\eta} 
= \frac{\eta}{2}\|\tilde{g}_t\|_2^2\,.
\end{equation*}
Using the previous inequality, we have: 
\begin{equation*}
\langle \tilde{g}_t, z_t - z_{\delta}^{\star} \rangle 
\leq \frac{D_R(z_{\delta}^{\star}||z_{t}) - D_R(z_{\delta}^{\star}||z_{t+1})}{\eta} + \frac{D_R(z_{\delta}^{\star}||z_{t+1}) - D_R(z_{\delta}^{\star}||y_{t+1})}{\eta} + \frac{\eta}{2}\|\tilde{g}_2\|_2^2\,.
\end{equation*}
By $G$-Lipschitzness of $z \mapsto D_R(z_{\delta}^{\star}||z)$, we have: 
\begin{equation*}
D_R(z_{\delta}^{\star}||z_{t+1}) - D_R(z_{\delta}^{\star}||y_{t+1}) 
\leq G \|z_{t+1} - y_{t+1}\|\,,
\end{equation*}
which implies that: 
\begin{equation*}
\langle \tilde{g}_t, z_t - z_{\delta}^{\star} \rangle 
\leq \frac{D_R(z_{\delta}^{\star}||z_{t}) - D_R(z_{\delta}^{\star}||z_{t+1})}{\eta} + \frac{G}{\eta} \|z_{t+1} - y_{t+1}\| + \frac{\eta}{2}\|\tilde{g}_2\|_2^2\,.
\end{equation*}
Summing up the above inequality, the first term telescopes and we obtain:  
\begin{equation*}
\sum_{t=1}^T \langle \tilde{g}_t, z_t - z_{\delta}^{\star} \rangle 
\leq \frac{D_R(z_{\delta}^{\star}||z_1)}{\eta} + \frac{\eta}{2} \sum_{t=1}^T\|\tilde{g}_t\|_2^2 + \frac{G}{\eta} \sum_{t=1}^T \|z_{t+1} - y_{t+1}\|_2\,,
\end{equation*}
which concludes the proof upon recalling the notation $r_{t+1} = z_{t+1} - y_{t+1}\,.$
As for the second bound, it follows immediately from using $\|\tilde{g}_t\| \leq \tilde{G}_F$ and $\|r_{t+1}\| \leq C_{\text{app}}$\,.
\end{proof}

In the next lemma, we show how to control the error sequence $(r_t)$ in Lemma~\ref{lem:perturbed-omd-reg-random-lin-loss} by providing a precise constant $\hat{C}_{\eta}$ which quantifies the approximation error in the bandit setting, as a corollary of Proposition~\ref{prop:dist-minimizers-eta2}. 

\begin{tcolorbox}[colframe=white!, top=2pt,left=2pt,right=2pt,bottom=2pt]
\begin{lemma}[Gradient estimators norm bounds]
\label{lem:grad-estimator-bounds}
Suppose that $\ell_t$ is uniformly bounded by $M$, i.e. $|\ell_t(x)| \leq M, \forall x \in \X$, and the inverse Jacobian of $q$ is bounded by $G$ at any point $x \in \X$ (as in Assumption~\ref{as:reg-R-reparam-q-smoothness}). Then for any $t \geq 1$, we have: 
\begin{align}
\|g_t\| &\leq \hat{G}_F := \frac{dM}{\delta}\,, \quad \|\tilde{g}_t\| \leq \tilde{G}_F := \frac{dM G}{\delta}\,.
\end{align}
\end{lemma}
\end{tcolorbox}

\begin{proof}
Recall from BOGD (Algorithm~\ref{alg4}) that $g_t = \frac{d}{\delta} \ell_t(\hat{x}_t) \zeta_t$. Then, since $\|\zeta_t\| = 1$, we have: 
\begin{equation*}
\|g_t\| = \frac{d}{\delta} |\ell_t(\hat{x}_t)| \leq \frac{dM}{\delta} = \hat{G}_F\,.
\end{equation*}
Then, since $\tilde{g}_t = J_q(x_t)^{-\top} g_t$, we have: 
\begin{equation*}
\|\tilde{g}_t\| \leq G \|g_t\| \leq \frac{dMG}{\delta} = \tilde{G}_F\,.
\end{equation*}
\end{proof}

\begin{tcolorbox}[colframe=white!, top=2pt,left=2pt,right=2pt,bottom=2pt]
\begin{corollary}[Proposition~\ref{prop:dist-minimizers-eta2} in the bandit setting]
\label{cor:approx-lemma-bandit}
In Proposition~\ref{prop:dist-minimizers-eta2}, let $g_t$ be the gradient estimator used in BOGD and recall that $\tilde{g}_t = J_q(x_t)^{-\top} g_t$ in \eqref{eq:omd}-\eqref{eq:ogd}, then the approximation inequality \eqref{eq:approx-ineq} holds with $\hat{G}_F = dM/\delta$ as defined in Lemma~\ref{lem:grad-estimator-bounds}, i.e., for any $t \geq 1$, if $y_t = q(x_t)$,  
\begin{equation*}
\|y_{t+1}-q(x_{t+1})\|_2 \leq \hat{C}_{\eta} :=  G^5 (5\hat{G}_F^2 + \hat{G}_F^3 \eta) \eta^2 \,.
\end{equation*}
\end{corollary}
\end{tcolorbox}

\noindent\textbf{End of Proof of Theorem~\ref{thm:bandit}.} We are ready to conclude the proof of the theorem using the results established so far. First, it follows from Lemma~\ref{lem:expected-reg-decomp} that: 
\begin{equation}
\label{eq:final-reg-decomp}
\mathbb{E}[R_T] \leq R_T^{(1)} + R_T^{(2)} + R_T^{(3)} + R_T^{(4)}\,.
\end{equation}
We have proved a bound for each one of the error terms in the above bound: 
\begin{enumerate}[label=(\roman*)]
    \item By Lemma~\ref{lem:domain-shrinkage-error}, we have $R_T^{(2)} \leq \delta G_FDT\,.$
    \item By Lemma~\ref{lem:grad-estimation-bias}, we get: $|R_T^{(3)}| \leq \delta G^2 H DT\,.$ 
    \item By Lemma~\ref{lem:smoothing-error}, we obtain: $| R_T^{(4)}| \leq \delta G_FT\,.$
    \item Lemma~\ref{lem:perturbed-omd-reg-random-lin-loss}-\eqref{eq:R1T-bandit} together with Corollary~\ref{cor:approx-lemma-bandit} yield: 
    \begin{equation*}
    R_T^{(1)} \leq \frac{D_1}{\eta} + \frac{\eta \tilde{G}_F^2T}{2} + \frac{\hat{C}_{\eta} G T}{\eta} \leq \frac{D_1}{\eta} + (6 \hat{G}_F^2 + \hat{G}_F^3 \eta) G^6 \eta T\,,
    \end{equation*}
    where the last inequality follows from using $\hat{C}_{\eta} =  G^5 (5\hat{G}_F^2 + \hat{G}_F^3 \eta) \eta^2$ (Corollary~\ref{cor:approx-lemma-bandit}) and upper bounding the resulting constant under the assumptions $G > 1$ (which is without loss of generality), similarly to the deterministic setting. Note here that $\hat{G}_F = \frac{dM}{\delta}$ as shown in Lemma~\ref{lem:grad-estimator-bounds}.   
\end{enumerate}
Combining all the above bounds in \eqref{eq:final-reg-decomp}, we obtain: 
\begin{equation}
\label{eq:regret-upper-bound-proof}
\mathbb{E}[R_T] \leq K_1 \delta T  + \frac{D_1}{\eta} + K_2 \frac{\eta T}{\delta^2} + K_3 \frac{\eta^2 T}{\delta^3}\,, 
\end{equation}
where $K_1 := G_FD + G^2 HD + G_F, K_2 := 6 d^2 G^6 M^2, K_3 = d^3 G^6 M^3\,.$

We now derive a choice of $(\eta, \delta)$ which gives an optimal dependence on the time horizon $T$. A lower bound on the upper regret bound (in $\eta, \delta$) is given by dropping the last positive error term $\frac{K_3 \eta^2 T}{\delta^3}$ (which as we will see will contribute to a faster regret rate and will hence not be the leading term). Define then for any $(\eta, \delta)$, 
\begin{equation*}
f(\eta, \delta) := K_1 \delta T  + \frac{D_1}{\eta} + K_2 \frac{\eta T}{\delta^2}\,, 
\end{equation*}
which consists of the three first terms in the upper bound \eqref{eq:regret-upper-bound-proof}. We now optimize this function w.r.t $(\eta, \delta)$. By the first order optimality conditions, the optimal pair $(\eta_{\star}, \delta_{\star})$ minimizing $f(\eta, \delta)$ satisfies: 
\begin{equation*}
\frac{\partial f}{\partial \eta} (\eta_{\star}, \delta_{\star}) = - \frac{D_1}{\eta^2} + \frac{K_2 T}{\delta_{\star}^2} = 0\,,\quad
\frac{\partial f}{\partial \delta} (\eta_{\star}, \delta_{\star}) 
= K_1 T - \frac{2 K_2 \eta_{\star} T}{\delta_{\star}^3}\,.
\end{equation*}
Rewriting these conditions yields: 
\begin{equation*}
\eta_{\star} = \sqrt{\frac{D_1 \delta_{\star}^2}{K_2 T}}\,, \quad \eta_{\star} = \frac{K_1 \delta_{\star}^3}{2 K_2}\,.
\end{equation*}
Solving for $\delta_{\star}$ as a function of $T$ (eliminating $\eta$) and then plugging back $\delta_{\star}$ in the first expression of $\eta_{\star}$ yields: 
\begin{equation*}
\delta_{\star} = \left(\frac{4 K_2 D_1}{K_1^2 T}\right)^{1/4}\,, \quad \eta_{\star} = \left(\frac{4 D_1^3}{K_1^2 K_2 T^3}\right)^{1/4}\,.
\end{equation*}
Now, plugging back these values of $\eta_{\star}$ and $\delta_{\star}$ in \eqref{eq:regret-upper-bound-proof}, we obtain: 
\begin{equation*}
\mathbb{E}[R_T] \leq (K_1^2 K_2D_1)^{1/4} T^{3/4} + \left(\frac{K_1^2 D_1^3 K_3^4}{4 K_2^5}\right) T^{1/4}\,.
\end{equation*}
Recalling that $K_1 = G_FD + G^2 HD + G_F, K_2 = 6 d^2 G^6 M^2, K_3 = d^3 G^6 M^3$, we bound each one of the terms in the bound above to obtain:
\begin{align}
\mathbb{E}[R_T] &\leq 
\sqrt{3 dM G^3 (G_FD + G^2 HD + G_F) D_1^{1/2}} T^{3/4}
+ \sqrt{dM G^3 (G_FD + G^2 HD + G_F) D_1^{3/2}} T^{1/4}\nonumber\\
&\leq \sqrt{dM G^3 (G_FD + G^2 HD + G_F) (3 D_1^{1/2} + D_1^{3/2})} T^{3/4} \,,
\end{align}
which concludes the proof of Theorem~\ref{thm:bandit}.


\end{document}
